\begin{oxlemma}{Preservation}{lemma:preservation-app}
  If $\oxtypjudge{\oxglobalctx}{\oxemptyctx}{\oxkontctx}{\oxvarctx}{\oxexpr}{\oxsitype_1}{\oxvarctx_f} and
  \oxstorevalidity{\oxglobalctx}{\oxvarctx}{\oxstore}$ and
  $\oxwfkontctx{\oxglobalctx}{\oxvarctx}{\overline{\oxvalue}}{\oxkontctx}$ and
  $\oxreduce{\oxglobalctx}{\oxstore}{\oxexpr}{\oxstore^\prime}{\oxexpr^\prime}$, then there exists
  $\oxvarctx_i$ such that
  $\oxstorevalidity{\oxglobalctx}{\oxvarctx_i}{\oxstore^\prime}$ and
  $\oxwfkontctx{\oxglobalctx}{\oxvarctx_i}{\overline{\oxvalue}}{\oxkontctx}$ and
  $\oxtypjudge{\oxglobalctx}{\oxemptyctx}{\oxkontctx}{\oxvarctx_i}{\oxexpr^\prime}
  {\oxsitype_2}{\oxvarctx_f^\prime}$ and
  $\oxtunify[\oxcombine]{\oxemptyctx}{\oxkontctx}{\oxvarctx_f^\prime}
  {\oxsitype_2}{\oxsitype_1}{\oxvarctx_s}$ and there exists $\oxvarctx_o$ such
  that $\oxvarctx_f = \oxvarctx_s \oxintersect \oxvarctx_o$.
\end{oxlemma}

\paragraph{Proof} We proceed by induction on the derivation
\oxtypjudge{\oxglobalctx}{\oxemptyctx}{\oxkontctx}{\oxvarctx}{\oxexpr}{\oxtype}{\oxvarctx_f} \\

\oxpreservationcase{T-Move}{\TMove}{\oxplace}{\EMove}{
  \oxtupdate{\oxvarctx}{\oxplace}{\oxsidtype}
}

\begin{preservation}
  %% new store is valid
  \oxpresline{
    \oxstorevalidity{\oxglobalctx}{
      \oxtupdate{\oxvarctx}{\oxplace}{\oxsidtype}
    }{
      \oxvupdate{\oxstore}{\oxplace}{\oxuninit}
    }
  }{
    Applying \Lemma{lemma:stack-validity-related-envs} to
    $\oxsubtypectx{\oxglobalctx}{\oxemptyctx}{\oxvarctx}{
      \oxtupdate{\oxvarctx}{\oxplace}{\oxsidtype}
    }$ (immediate by \oxname{R-Env}) and
    $\oxstorevalidity{\oxglobalctx}{\oxvarctx}{\oxstore}$ (from premise) gives us
    $\oxstorevalidity{\oxglobalctx}{
      \oxtupdate{\oxvarctx}{\oxplace}{\oxsidtype}
    }{\oxstore}$. Then, since we know
    $\oxtypjudge{\oxglobalctx}{\oxemptyctx}{\oxkontctx}{
      \oxtupdate{\oxvarctx}{\oxplace}{\oxsidtype}
    }{\oxuninit}{\oxsidtype}{
      \oxtupdate{\oxvarctx}{\oxplace}{\oxsidtype}
    }$ (by \oxname{T-Dead}), we can conclude
    $\oxstorevalidity{\oxglobalctx}{
      \oxtupdate{\oxvarctx}{\oxplace}{\oxsidtype}
    }{
      \oxvupdate{\oxstore}{\oxplace}{\oxuninit}
    }$.
  }

  %% kontctx values are still valid
  \oxpresline{
    \oxwfkontctx{\oxglobalctx}{
      \oxtupdate{\oxvarctx}{\oxplace}{\oxsidtype}
    }{\overline{\oxvalue}}{\oxkontctx}
  }{
    Inverting \oxname{WF-Temporaries} gives us $\forall i \in 1 \oxdots n. \;
    \oxtypjudge{\oxglobalctx}{\oxemptyctx}{\oxsitype_1 \oxdotsc \oxsitype_{i-1}}{\oxvarctx}{
      \oxvalue_i
    }{\oxsitype_i}{\oxvarctx}$. By \oxname{R-Env}, we have that
    \oxsubtypectx{\oxglobalctx}{\oxemptyctx}{\oxvarctx}{
      \oxtupdate{\oxvarctx}{\oxplace}{\oxsidtype}
    }. Thus, we can apply \Lemma{lemma:value-typing-related-envs} to each of the
    typing judgments and then apply \oxname{WF-Temporaries} to get
    $\oxwfkontctx{\oxglobalctx}{
      \oxtupdate{\oxvarctx}{\oxplace}{\oxsidtype}
    }{\overline{\oxvalue}}{\oxkontctx}$.
  }

  %% new expression is well-typed
  \oxpresline{
    \oxtypjudge{\oxglobalctx}{\oxemptyctx}{\oxkontctx}{
      \oxtupdate{\oxvarctx}{\oxplace}{\oxsidtype}
    }{\oxvalue}{\oxsitype}{
      \oxtupdate{\oxvarctx}{\oxplace}{\oxsidtype}
    }
  }{
    Applying \Lemma{lemma:place-exprs-reduce} to
    $\oxmusafety{\oxemptyctx}{\oxvarctx}{\oxmut}{\oxplace}{
      \oxset{\oxloanpkg{\oxmut}{\oxplace}}
    }$, $\oxcomputetynoprov{\oxemptyctx}{\oxvarctx}{\oxmut}{\oxplace}{\oxsitype}$
    (immediate by \oxname{TC-Place} with $\oxlookup{\oxvarctx}{\oxplace}{
      \oxsitype
    }$), and $\oxstorevalidity{\oxglobalctx}{\oxvarctx}{\oxstore}$ gives us
    $\oxtypjudge{\oxglobalctx}{\oxemptyctx}{\oxkontctx}{\oxvarctx}{\oxvalue}
    {\oxsitype}{\oxvarctx}$. Then, by applying
    \Lemma{lemma:value-typing-related-envs} with
    $\oxsubtypectx{\oxglobalctx}{\oxemptyctx}
    {\oxvarctx}{\oxtupdate{\oxvarctx}{\oxplace}{\oxsidtype}}$ (immediate by
    \oxname{R-Env}), we can conclude $\oxtypjudge{\oxglobalctx}{\oxemptyctx}{\oxkontctx}{
      \oxtupdate{\oxvarctx}{\oxplace}{\oxsidtype}
    }{\oxvalue}{\oxsitype}{
      \oxtupdate{\oxvarctx}{\oxplace}{\oxsidtype}
    }$.
  }

  %% new type is a subtype of the old type
  \oxpresline{
    \oxtunify[\oxcombine]{\oxemptyctx}{\oxkontctx}{
      \oxtupdate{\oxvarctx}{\oxplace}{\oxsidtype}
    }{\oxsitype}{\oxsitype}{
      \oxtupdate{\oxvarctx}{\oxplace}{\oxsidtype}
    }
  }{
    Immediate by \oxname{RR-Refl}.
  }

  %% new loan context is a sub-context of old loan context
  \oxpresline{
    \exists \oxvarctx_o. \oxtupdate{\oxvarctx}{\oxplace}{\oxsidtype} \oxintersect \oxvarctx_o
    = \oxtupdate{\oxvarctx}{\oxplace}{\oxsidtype}
  }{
    $\oxvarctx_o = \oxtupdate{\oxvarctx}{\oxplace}{\oxsidtype}$
  }
\end{preservation}

\oxpreservationcase{T-Copy}{\TCopy}{\oxplaceexpr}{\ECopy}{
  \oxvarctx
}

\begin{preservation}
  %% new store is valid
  \oxpresline{
    \oxstorevalidity{\oxglobalctx}{\oxvarctx}{\oxstore}
  }{
    Immediate from our premise.
  }

  %% kontctx values are still valid
  \oxpresline{
    \oxwfkontctx{\oxglobalctx}{\oxvarctx}{\overline{\oxvalue}}{\oxkontctx}
  }{
    Immediate from our premise.
  }

  %% new expression is well-typed
  \oxpresline{
    \oxtypjudge{\oxglobalctx}{\oxemptyctx}{\oxkontctx}{\oxvarctx}{\oxvalue}
    {\oxsitype}{\oxvarctx}
  }{
    Applying \Lemma{lemma:place-exprs-reduce} to
    $\oxmusafety{\oxemptyctx}{\oxvarctx}{\oximm}{\oxplaceexpr}{
      \oxset{\overline{\oxloan}}
    }$, $\oxcomputetynoprov{\oxemptyctx}{\oxvarctx}{\oximm}{\oxplaceexpr}
    {\oxsitype}$, and $\oxstorevalidity{\oxglobalctx}{\oxvarctx}{\oxstore}$
    gives us $\oxtypjudge{\oxglobalctx}{\oxemptyctx}{\oxkontctx}{\oxvarctx}{\oxvalue}
    {\oxsitype}{\oxvarctx}$.
  }

  %% new type is a subtype of the old type
  \oxpresline{
    \oxtunify[\oxcombine]{\oxemptyctx}{\oxkontctx}{\oxvarctx}
    {\oxsitype}{\oxsitype}{\oxvarctx}
  }{
    Immediate by \oxname{RR-Refl}.
  }

  %% new loan context is a sub-context of the old loan context
  \oxpresline{
    \exists \oxvarctx_o. \oxvarctx \oxintersect \oxvarctx_o = \oxvarctx
  }{
    $\oxvarctx_o = \oxvarctx$
  }
\end{preservation}

\oxpreservationcase{T-Borrow}{\TBorrow}{
  \oxref{\oxprov}{\oxmuta}{\oxplaceexpr}
}{\EBorrow}{
  \oxupdate{\oxvarctx}{
    \oxloanctxentry{\oxcprov}{
      \oxset{\overline{\oxloan}}
    }
  }
}

\begin{preservation}
  %% new store is valid
  \oxpresline{
    \oxstorevalidity{\oxglobalctx}{
      \oxupdate{\oxvarctx}{
        \oxloanctxentry{\oxcprov}{
          \oxset{\overline{\oxloan}}
        }
      }
    }{\oxstore}
  }{
    Apply \Lemma{lemma:store-validity-loan-update} to
    $\oxstorevalidity{\oxglobalctx}{\oxvarctx}{\oxstore}$ (from our premise)
    and $\oxctxswellformed{\oxglobalctx}{\oxemptyctx}{
      \oxupdate{\oxvarctx}{
        \oxloanctxentry{\oxcprov}{
          \oxset{\overline{\oxloan}}
        }
      }
    }{\oxkontctx}$ (from our premise)
    gives us $\oxstorevalidity{\oxglobalctx}{
      \oxupdate{\oxvarctx}{
        \oxloanctxentry{\oxcprov}{
          \oxset{\overline{\oxloan}}
        }
      }
    }{\oxstore}$.
  }

  %% kontctx values are still valid
  \oxpresline{
    \oxwfkontctx{\oxglobalctx}{
      \oxupdate{\oxvarctx}{
        \oxloanctxentry{\oxcprov}{
          \oxset{\overline{\oxloan}}
        }
      }
    }{\overline{\oxvalue}}{\oxkontctx}
  }{
    Inverting \oxname{WF-Temporaries} gives us $\forall i \in 1 \oxdots n. \;
    \oxtypjudge{\oxglobalctx}{\oxemptyctx}{\oxsitype_1 \oxdotsc \oxsitype_{i-1}}{\oxvarctx}{
      \oxvalue_i
    }{\oxsitype_i}{\oxvarctx}$. We can then apply
    \Lemma{lemma:values-well-typed-loan-update} to each of the typing judgments
    along with
    $\oxmusafety{\oxemptyctx}{\oxkontctx}{\oxvarctx}{\oxmuta}{\oxplaceexpr}{
      \oxset{\overline{\oxloan}}
    }$ and $\oxnotinclosure{\oxkontctx}{\oxvarctx}{\oxcprov}$
    and $\oxlookup{\oxvarctx}{\oxcprov}{\emptyset}$ (all from the premise of
    \oxname{T-Borrow}) to get $\forall i \in 1 \oxdots n. \;
    \oxtypjudge{\oxglobalctx}{\oxemptyctx}{\oxsitype_1 \oxdotsc \oxsitype_{i-1}}{
      \oxupdate{\oxvarctx}{
        \oxloanctxentry{\oxcprov}{
          \oxset{\overline{\oxloan}}
        }
      }
    }{\oxvalue_i}{\oxsitype_i}{
      \oxupdate{\oxvarctx}{
        \oxloanctxentry{\oxcprov}{
          \oxset{\overline{\oxloan}}
        }
      }
    }$. We can then apply \oxname{WF-Temporaries} to get
    $\oxwfkontctx{\oxglobalctx}{
      \oxupdate{\oxvarctx}{
        \oxloanctxentry{\oxcprov}{
          \oxset{\overline{\oxloan}}
        }
      }
    }{\overline{\oxvalue}}{\oxkontctx}$.
  }

  %% new expression is well-typed
  \oxpresline{
    \oxtypjudge{\oxglobalctx}{\oxemptyctx}{\oxkontctx}{\oxvarctx_i}{
      \oxptr{\oxreferent}
    }{\oxtref{\oxcprov}{\oxmuta}{\oxxitype}}{\oxvarctx_i}
  }{
    Applying \Lemma{lemma:norm-for-valid-referents} to
    $\oxstorevalidity{\oxglobalctx}{\oxvarctx}{\oxstore}$, and
    $\oxnorm{\oxstore}{\oxplaceexpr}{\oxreferent}{\_}{\_}$ gives us
     $\oxreferentvalidity{\oxglobalctx}{\oxvarctx}{
      \oxreferentctx[\oxplace]
    }{\oxxitype}$. Then, note that referent well-formedness does not
    depend on the contents of loan sets. This means we can also conclude
    $\oxreferentvalidity{\oxglobalctx}{
      \oxupdate{\oxvarctx}{
        \oxloanctxentry{\oxcprov}{
          \oxset{\overline{\oxloan}}
        }
      }
    }{
      \oxreferentctx[\oxplace]
    }{\oxxitype}$.

    Applying \Lemma{lemma:norm-place-prefix-in-loans} to
    $\oxstorevalidity{\oxglobalctx}{\oxvarctx}{\oxstore}$,
    $\oxnorm{\oxstore}{\oxplaceexpr}{\oxreferent}{\_}{\_}$, and
    $\oxmusafety{\oxemptyctx}{\oxvarctx}{\oxmuta}{\oxplaceexpr}{
      \oxset{\overline{\oxloan}}
    }$ gives us $\oxreferent = \oxreferentctx[\oxplace]$ and
    $\oxloanpkg{\oxmuta}{\oxplace} \in \oxset{\overline{\oxloan}}$.

    Finally, we can apply \oxname{T-Pointer} to the two facts above to get
    $\oxtypjudge{\oxglobalctx}{\oxemptyctx}{\oxkontctx}{
      \oxupdate{\oxvarctx}{
        \oxloanctxentry{\oxcprov}{
          \oxset{\overline{\oxloan}}
        }
      }
    }{
      \oxptr{\oxreferent}
    }{\oxtref{\oxcprov}{\oxmuta}{\oxxitype}}{
      \oxupdate{\oxvarctx}{
        \oxloanctxentry{\oxcprov}{
          \oxset{\overline{\oxloan}}
        }
      }
    }$.
  }

  %% new type is a subtype of the old type
  \oxpresline{
    \oxtunify[\oxcombine]{\oxemptyctx}{\oxkontctx}{\oxvarctx_i}{
      \oxtref{\oxcprov}{\oxmuta}{\oxxitype}
    }{
      \oxtref{\oxcprov}{\oxmuta}{\oxxitype}
    }{\oxvarctx_i}
  }{
    Immediate by \oxname{RR-Refl}.
  }

  %% new loan context is a sub-context of the old loan context
  \oxpresline{
    \exists \oxvarctx_o. \oxvarctx_i \oxintersect \oxvarctx_o = \oxvarctx_i
  }{
    $\oxvarctx_o = \oxvarctx_i$
  }
\end{preservation}

\oxpreservationcasemulti{T-BorrowIndex}{\TBorrowIndex}{
  \oxref{\oxprov}{\oxmuta}{\oxindex{\oxplaceexpr}{\oxexpr_i}}
}{
  \EBorrowIndex \and \EBorrowIndexOOB
}

\oxpreservationsubcaseheading{E-BorrowIndex}{
  \oxupdate{\oxvarctx^\prime}{
    \oxloanctxentry{\oxcprov}{
      \oxset{\overline{\oxloan}}
    }
  }
}

\begin{preservation}
  %% new store is valid
  \oxpresline{
    \oxstorevalidity{\oxglobalctx}{
      \oxupdate{\oxvarctx^\prime}{
        \oxloanctxentry{\oxcprov}{
          \oxset{\overline{\oxloan}}
        }
      }
    }{\oxstore}
  }{
    Applying \Lemma{lemma:values-dont-change} to the typing derivation (from
    \oxname{T-BorrowIndex}) for
    $\oxexpr$ (which we know is a value from \oxname{E-BorrowIndex}) gives us
    $\oxsubtypectx{\oxglobalctx}{\oxemptyctx}{\oxvarctx}{\oxvarctx^\prime}$.

    Then, applying \Lemma{lemma:stack-validity-related-envs} to
    $\oxsubtypectx{\oxglobalctx}{\oxemptyctx}{\oxvarctx}{\oxvarctx^\prime}$ and
    $\oxstorevalidity{\oxglobalctx}{\oxvarctx}{\oxstore}$ (from premise) gives us
    $\oxstorevalidity{\oxglobalctx}{\oxvarctx^\prime}{\oxstore}$.

    Finally, applying \Lemma{lemma:store-validity-loan-update} to
    $\oxstorevalidity{\oxglobalctx}{\oxvarctx^\prime}{\oxstore}$ gives us
    $\oxstorevalidity{\oxglobalctx}{
      \oxupdate{\oxvarctx^\prime}{
        \oxloanctxentry{\oxcprov}{
          \oxset{\overline{\oxloan}}
        }
      }
    }{\oxstore}$.
  }

  %% kontctx values are still valid
  \oxpresline{
    \oxwfkontctx{\oxglobalctx}{
      \oxupdate{\oxvarctx^\prime}{
        \oxloanctxentry{\oxcprov}{
          \oxset{\overline{\oxloan}}
        }
      }
    }{\overline{\oxvalue}}{\oxkontctx}
  }{
    Inverting \oxname{WF-Temporaries} gives us $\forall i \in 1 \oxdots n. \;
    \oxtypjudge{\oxglobalctx}{\oxemptyctx}{\oxsitype_1 \oxdotsc \oxsitype_{i-1}}{\oxvarctx}{
      \oxvalue_i
    }{\oxsitype_i}{\oxvarctx}$.

    Applying \Lemma{lemma:values-dont-change} to the typing derivation (from
    \oxname{T-BorrowIndex}) for
    $\oxexpr$ (which we know is a value from \oxname{E-BorrowIndex}) gives us
    $\oxsubtypectx{\oxglobalctx}{\oxemptyctx}{\oxvarctx}{\oxvarctx^\prime}$.

    Thus, we can apply \Lemma{lemma:value-typing-related-envs} to each of the
    typing judgments to get $\forall i \in 1 \oxdots n. \;
    \oxtypjudge{\oxglobalctx}{\oxemptyctx}{\oxemptyctx}{\oxvarctx^\prime}{
      \oxvalue_i
    }{\oxsitype_i}{\oxvarctx^\prime}$.

    We can then apply \Lemma{lemma:values-well-typed-loan-update} to each of the
    typing judgments along with
    $\oxmusafety{\oxemptyctx}{\oxkontctx}{\oxvarctx^\prime}{\oxmuta}{\oxplaceexpr}{
      \oxset{\overline{\oxloan}}
    }$ and $\oxnotinclosure{\oxkontctx}{\oxvarctx^\prime}{\oxcprov}$
    and $\oxlookup{\oxvarctx^\prime}{\oxcprov}{\emptyset}$ (all from the premise
    of \oxname{T-BorrowIndex}) to get $\forall i \in 1 \oxdots n. \;
    \oxtypjudge{\oxglobalctx}{\oxemptyctx}{\oxemptyctx}{
      \oxupdate{\oxvarctx^\prime}{
        \oxloanctxentry{\oxcprov}{
          \oxset{\overline{\oxloan}}
        }
      }
    }{\oxvalue_i}{\oxsitype_i}{
      \oxupdate{\oxvarctx^\prime}{
        \oxloanctxentry{\oxcprov}{
          \oxset{\overline{\oxloan}}
        }
      }
    }$. We can then apply \oxname{WF-Temporaries} to get
    $\oxwfkontctx{\oxglobalctx}{
      \oxupdate{\oxvarctx^\prime}{
        \oxloanctxentry{\oxcprov}{
          \oxset{\overline{\oxloan}}
        }
      }
    }{\overline{\oxvalue}}{\oxkontctx}$.
  }

  %% new expression is well-typed
  \oxpresline{
    \oxtypjudge{\oxglobalctx}{\oxemptyctx}{\oxkontctx}{\oxvarctx_i}{
      \oxindexptr{\oxreferent}{\oxnum_i}
    }{\oxtref{\oxcprov}{\oxmuta}{\oxsitype}}{\oxvarctx_i}
  }{
    Applying \Lemma{lemma:values-dont-change} to the typing derivation (from
    \oxname{T-BorrowIndex}) for
    $\oxexpr$ (which we know is a value from \oxname{E-BorrowIndex}) gives us
    $\oxsubtypectx{\oxglobalctx}{\oxemptyctx}{\oxvarctx}{\oxvarctx^\prime}$.

    Then, applying \Lemma{lemma:stack-validity-related-envs} to
    $\oxsubtypectx{\oxglobalctx}{\oxemptyctx}{\oxvarctx}{\oxvarctx^\prime}$ and
    $\oxstorevalidity{\oxglobalctx}{\oxvarctx}{\oxstore}$ (from premise) gives us
    $\oxstorevalidity{\oxglobalctx}{\oxvarctx^\prime}{\oxstore}$.

    Applying \Lemma{lemma:norm-for-valid-referents} to
    $\oxstorevalidity{\oxglobalctx}{\oxvarctx^\prime}{\oxstore}$, and
    $\oxnorm{\oxstore}{\oxplaceexpr}{\oxreferent}{\_}{
      \oxarr{\oxvalue_0 \oxdotsc \oxvalue_n}
    }$ gives us $\oxreferentvalidity{\oxglobalctx}{\oxvarctx^\prime}{
      \oxreferentctx[\oxplace]
    }{\oxxitype}$. Then, note that referent well-formedness does not
    depend on the contents of loan sets. This means we can also conclude
    $\oxreferentvalidity{\oxglobalctx}{
      \oxupdate{\oxvarctx^\prime}{
        \oxloanctxentry{\oxcprov}{
          \oxset{\overline{\oxloan}}
        }
      }
    }{
      \oxreferentctx[\oxplace]
    }{\oxxitype}$. We can then apply \oxname{WF-RefIndexArray} or
    \oxname{WF-RefIndexSlice} to get $\oxreferentvalidity{\oxglobalctx}{
      \oxupdate{\oxvarctx^\prime}{
        \oxloanctxentry{\oxcprov}{
          \oxset{\overline{\oxloan}}
        }
      }
    }{
      \oxindex{\oxreferentctx[\oxplace]}{\oxnum_i}
    }{\oxxitype}$.

    Applying \Lemma{lemma:norm-place-prefix-in-loans} to
    $\oxstorevalidity{\oxglobalctx}{\oxvarctx^\prime}{\oxstore}$,
    $\oxnorm{\oxstore}{\oxplaceexpr}{\oxreferent}{\_}{
      \oxarr{\oxvalue_0 \oxdotsc \oxvalue_n}
    }$, and
    $\oxmusafety{\oxemptyctx}{\oxvarctx^\prime}{\oxmuta}{\oxplaceexpr}{
      \oxset{\overline{\oxloan}}
    }$ gives us $\oxreferent = \oxreferent[\oxplace]$ and
    $\oxloanpkg{\oxmuta}{\oxplace} \in \oxset{\overline{\oxloan}}$.

    Finally, we can apply \oxname{T-Pointer} to the two facts above to get
    $\oxtypjudge{\oxglobalctx}{\oxemptyctx}{\oxkontctx}{
      \oxupdate{\oxvarctx}{
        \oxloanctxentry{\oxcprov}{
          \oxset{\overline{\oxloan}}
        }
      }
    }{
      \oxindexptr{\oxreferent}{\oxnum_i}
    }{\oxtref{\oxcprov}{\oxmuta}{\oxxitype}}{
      \oxupdate{\oxvarctx}{
        \oxloanctxentry{\oxcprov}{
          \oxset{\overline{\oxloan}}
        }
      }
    }$.
  }

  %% new type is a subtype of the old type
  \oxpresline{
    \oxtunify[\oxcombine]{\oxemptyctx}{\oxkontctx}{\oxvarctx_i}{
      \oxtref{\oxcprov}{\oxmuta}{\oxsitype}
    }{
      \oxtref{\oxcprov}{\oxmuta}{\oxsitype}
    }{\oxvarctx_i}
  }{
    Immediate by \oxname{RR-Refl}.
  }

  %% new loan context is a sub-context of the old loan context
  \oxpresline{
    \exists \oxvarctx_o. \oxvarctx_i \oxintersect \oxvarctx_o = \oxvarctx_i
  }{
    $\oxvarctx_o = \oxvarctx_i$
  }
\end{preservation}

\oxpreservationsubcaseheading{E-BorrowIndexOOB}{
  \oxvarctx
}

\begin{preservation}
  %% new store is valid
  \oxpresline{
    \oxstorevalidity{\oxglobalctx}{\oxvarctx}{\oxstore}
  }{
    Immediate from our premise.
  }

  %% kontctx values are still valid
  \oxpresline{
    \oxwfkontctx{\oxglobalctx}{\oxvarctx}{\overline{\oxvalue}}{\oxkontctx}
  }{
    Immediate from our premise.
  }

  %% new expression is well-typed
  \oxpresline{
    \oxtypjudge{\oxglobalctx}{\oxemptyctx}{\oxkontctx}{\oxvarctx}
    {\oxkey{abort!}(\oxdots)}{\oxtref{\oxcprov}{\oxmuta}{\oxsitype}}{\oxvarctx}
  }{
    An \oxkey{abort!} expression is well-typed (at any type) via the rule
    \oxname{T-Abort}.
  }

  %% new type is a subtype of the old type
  \oxpresline{
    \oxtunify[\oxcombine]{\oxemptyctx}{\oxkontctx}{\oxvarctx}{
      \oxtref{\oxcprov}{\oxmuta}{\oxsitype}
    }{
      \oxtref{\oxcprov}{\oxmuta}{\oxsitype}
    }{\oxvarctx}
  }{
    Immediate by \oxname{RR-Refl}.
  }

  %% new loan context is a sub-context of the old loan context
  \oxpresline{
    \exists \oxvarctx_o. \oxvarctx_i \oxintersect \oxvarctx_o = \oxvarctx_i
  }{
    $\oxvarctx_o = \oxvarctx_i$
  }
\end{preservation}

\oxpreservationcasemulti{T-BorrowSlice}{\TBorrowSlice}{
  \oxref{\oxcprov}{\oxmuta}{\oxslice{\oxplaceexpr}{\oxexpr_1}{\oxexpr_2}}
}{
  \EBorrowSlice \and \EBorrowSliceOOB
}

\oxpreservationsubcaseheading{E-BorrowSlice}{
  \oxupdate{\oxvarctx_2}{
    \oxloanctxentry{\oxcprov}{
      \oxset{\overline{\oxloan}}
    }
  }
}

\begin{preservation}
  %% new store is valid
  \oxpresline{
    \oxstorevalidity{\oxglobalctx}{
      \oxupdate{\oxvarctx_2}{
        \oxloanctxentry{\oxcprov}{
          \oxset{\overline{\oxloan}}
        }
      }
    }{\oxstore}
  }{
    Applying \Lemma{lemma:values-dont-change} to the typing derivation (from
    \oxname{T-BorrowSlice}) for $\oxexpr_1$ (which we know is a value from
    \oxname{E-BorrowSlice}) gives us
    $\oxsubtypectx{\oxglobalctx}{\oxemptyctx}{\oxvarctx}{\oxvarctx_1}$. Then,
    applying \Lemma{lemma:values-dont-change} to the typing derivation (from
    \oxname{T-BorrowSlice}) for  $\oxexpr_2$ (which we know is a value from
    \oxname{E-BorrowSlice}) gives us
    $\oxsubtypectx{\oxglobalctx}{\oxemptyctx}{\oxvarctx_1}{\oxvarctx_2}$. Then,
    by transitivity, we have
    $\oxsubtypectx{\oxglobalctx}{\oxemptyctx}{\oxvarctx}{\oxvarctx_2}$.

    Then, applying \Lemma{lemma:stack-validity-related-envs} to
    $\oxsubtypectx{\oxglobalctx}{\oxemptyctx}{\oxvarctx}{\oxvarctx_2}$ and
    $\oxstorevalidity{\oxglobalctx}{\oxvarctx}{\oxstore}$ (from premise) gives us
    $\oxstorevalidity{\oxglobalctx}{\oxvarctx_2}{\oxstore}$. Finally,
    applying \Lemma{lemma:store-validity-loan-update} to
    $\oxstorevalidity{\oxglobalctx}{\oxvarctx_2}{\oxstore}$ gives us
    $\oxstorevalidity{\oxglobalctx}{
      \oxupdate{\oxvarctx_2}{
        \oxloanctxentry{\oxcprov}{
          \oxset{\overline{\oxloan}}
        }
      }
    }{\oxstore}$.
  }

  %% kontctx values are still valid
  \oxpresline{
    \oxwfkontctx{\oxglobalctx}{
      \oxupdate{\oxvarctx_2}{
        \oxloanctxentry{\oxcprov}{
          \oxset{\overline{\oxloan}}
        }
      }
    }{\overline{\oxvalue}}{\oxkontctx}
  }{
    Inverting \oxname{WF-Temporaries} gives us $\forall i \in 1 \oxdots n. \;
    \oxtypjudge{\oxglobalctx}{\oxemptyctx}{\oxsitype_1 \oxdotsc \oxsitype_{i-1}}{\oxvarctx}{
      \oxvalue_i
    }{\oxsitype_i}{\oxvarctx}$.

    Applying \Lemma{lemma:values-dont-change} to the typing derivation (from
    \oxname{T-BorrowSlice}) for $\oxexpr_1$ (which we know is a value from
    \oxname{E-BorrowSlice}) gives us
    $\oxsubtypectx{\oxglobalctx}{\oxemptyctx}{\oxvarctx}{\oxvarctx_1}$. Then,
    applying \Lemma{lemma:values-dont-change} to the typing derivation (from
    \oxname{T-BorrowSlice}) for  $\oxexpr_2$ (which we know is a value from
    \oxname{E-BorrowSlice}) gives us
    $\oxsubtypectx{\oxglobalctx}{\oxemptyctx}{\oxvarctx_1}{\oxvarctx_2}$. Then,
    by transitivity, we have
    $\oxsubtypectx{\oxglobalctx}{\oxemptyctx}{\oxvarctx}{\oxvarctx_2}$.

    Then, we apply \Lemma{lemma:value-typing-related-envs} to each of the
    typing judgments to get $\forall i \in 1 \oxdots n. \;
    \oxtypjudge{\oxglobalctx}{\oxemptyctx}{\oxsitype_1 \oxdotsc \oxsitype_{i-1}}{\oxvarctx_2}{
      \oxvalue_i
    }{\oxsitype_i}{\oxvarctx_2}$.

    We can then apply \Lemma{lemma:values-well-typed-loan-update} to each of the
    typing judgments along with
    $\oxmusafety{\oxemptyctx}{\oxkontctx}{\oxvarctx_2}{\oxmuta}{\oxplaceexpr}{
      \oxset{\overline{\oxloan}}
    }$ and $\oxnotinclosure{\oxkontctx}{\oxvarctx_2}{\oxcprov}$
    and $\oxlookup{\oxvarctx_2}{\oxcprov}{\emptyset}$ (all from the premise
    of \oxname{T-BorrowSlice}) to get $\forall i \in 1 \oxdots n. \;
    \oxtypjudge{\oxglobalctx}{\oxemptyctx}{\oxsitype_1 \oxdotsc \oxsitype_{i-1}}{
      \oxupdate{\oxvarctx_2}{
        \oxloanctxentry{\oxcprov}{
          \oxset{\overline{\oxloan}}
        }
      }
    }{\oxvalue_i}{\oxsitype_i}{
      \oxupdate{\oxvarctx_2}{
        \oxloanctxentry{\oxcprov}{
          \oxset{\overline{\oxloan}}
        }
      }
    }$. We can then apply \oxname{WF-Temporaries} to get
    $\oxwfkontctx{\oxglobalctx}{
      \oxupdate{\oxvarctx_2}{
        \oxloanctxentry{\oxcprov}{
          \oxset{\overline{\oxloan}}
        }
      }
    }{\overline{\oxvalue}}{\oxkontctx}$.
  }

  %% new expression is well-typed
  \oxpresline{
    \oxtypjudge{\oxglobalctx}{\oxemptyctx}{\oxkontctx}{\oxvarctx_i}{
      \oxsliceptr{\oxreferent}{\oxnum_1}{\oxnum_2}
    }{
      \oxtref{\oxcprov}{\oxmuta}{\oxtslice{\oxsitype}}
    }{\oxvarctx_i}
  }{
    Applying \Lemma{lemma:values-dont-change} to the typing derivation (from
    \oxname{T-BorrowSlice}) for $\oxexpr_1$ (which we know is a value from
    \oxname{E-BorrowSlice}) gives us
    $\oxsubtypectx{\oxglobalctx}{\oxemptyctx}{\oxvarctx}{\oxvarctx_1}$. Then,
    applying \Lemma{lemma:values-dont-change} to the typing derivation (from
    \oxname{T-BorrowSlice}) for  $\oxexpr_2$ (which we know is a value from
    \oxname{E-BorrowSlice}) gives us
    $\oxsubtypectx{\oxglobalctx}{\oxemptyctx}{\oxvarctx_1}{\oxvarctx_2}$. Then,
    by transitivity, we have
    $\oxsubtypectx{\oxglobalctx}{\oxemptyctx}{\oxvarctx}{\oxvarctx_2}$.

    Then, applying \Lemma{lemma:stack-validity-related-envs} to
    $\oxsubtypectx{\oxglobalctx}{\oxemptyctx}{\oxvarctx}{\oxvarctx_2}$ and
    $\oxstorevalidity{\oxglobalctx}{\oxvarctx}{\oxstore}$ (from premise) gives us
    $\oxstorevalidity{\oxglobalctx}{\oxvarctx_2}{\oxstore}$.

    Applying \Lemma{lemma:norm-for-valid-referents} to
    $\oxstorevalidity{\oxglobalctx}{\oxvarctx_2}{\oxstore}$, and
    $\oxnorm{\oxstore}{\oxplaceexpr}{\oxreferent}{\_}{
      \oxarr{\oxvalue_0 \oxdotsc \oxvalue_n}
    }$ gives us $\oxreferentvalidity{\oxglobalctx}{\oxvarctx_2}{
      \oxreferentctx[\oxplace]
    }{\oxxitype}$. Then, note that referent well-formedness does not
    depend on the contents of loan sets. This means we can also conclude
    $\oxreferentvalidity{\oxglobalctx}{
      \oxupdate{\oxvarctx_2}{
        \oxloanctxentry{\oxcprov}{
          \oxset{\overline{\oxloan}}
        }
      }
    }{
      \oxreferentctx[\oxplace]
    }{\oxxitype}$. We can then apply \oxname{WF-RefSliceArray} or
    \oxname{WF-RefSliceSlice} to get $\oxreferentvalidity{\oxglobalctx}{
      \oxupdate{\oxvarctx^\prime}{
        \oxloanctxentry{\oxcprov}{
          \oxset{\overline{\oxloan}}
        }
      }
    }{
      \oxslice{\oxreferentctx[\oxplace]}{\oxnum_1}{\oxnum_2}
    }{\oxxitype}$.

    Applying \Lemma{lemma:norm-place-prefix-in-loans} to
    $\oxstorevalidity{\oxglobalctx}{\oxvarctx_2}{\oxstore}$,
    $\oxnorm{\oxstore}{\oxplaceexpr}{\oxreferent}{\_}{
      \oxarr{\oxvalue_0 \oxdotsc \oxvalue_n}
    }$, and
    $\oxmusafety{\oxemptyctx}{\oxvarctx_2}{\oxmuta}{\oxplaceexpr}{
      \oxset{\overline{\oxloan}}
    }$ gives us $\oxloanpkg{\oxmuta}{\oxplace} \in \oxset{\overline{\oxloan}}$.

    Finally, we can apply \oxname{T-Pointer} to the two facts above to get
    $\oxtypjudge{\oxglobalctx}{\oxemptyctx}{\oxkontctx}{
      \oxupdate{\oxvarctx}{
        \oxloanctxentry{\oxcprov}{
          \oxset{\overline{\oxloan}}
        }
      }
    }{
      \oxsliceptr{\oxreferent[\oxplace]}{\oxnum_1}{\oxnum_2}
    }{\oxtref{\oxcprov}{\oxmuta}{\oxxitype}}{
      \oxupdate{\oxvarctx}{
        \oxloanctxentry{\oxcprov}{
          \oxset{\overline{\oxloan}}
        }
      }
    }$.
  }

  %% new type is a subtype of the old type
  \oxpresline{
    \oxtunify[\oxcombine]{\oxemptyctx}{\oxkontctx}{\oxvarctx_i}{
      \oxtref{\oxcprov}{\oxmuta}{\oxtslice{\oxsitype}}
    }{
      \oxtref{\oxcprov}{\oxmuta}{\oxtslice{\oxsitype}}
    }{\oxvarctx_i}
  }{
    Immediate by \oxname{RR-Refl}.
  }

  %% new loan context is a sub-context of the old loan context
  \oxpresline{
    \exists \oxvarctx_o. \oxvarctx_i \oxintersect \oxvarctx_o = \oxvarctx_i
  }{
    $\oxvarctx_o = \oxvarctx_i$
  }
\end{preservation}

\oxpreservationsubcaseheading{E-BorrowSliceOOB}{
  \oxvarctx
}

\begin{preservation}
  %% new store is valid
  \oxpresline{
    \oxstorevalidity{\oxglobalctx}{\oxvarctx}{\oxstore}
  }{
    Immediate from our premise.
  }

  %% kontctx values are still valid
  \oxpresline{
    \oxwfkontctx{\oxglobalctx}{\oxvarctx}{\overline{\oxvalue}}{\oxkontctx}
  }{
    Immediate from our premise.
  }

  %% new expression is well-typed
  \oxpresline{
    \oxtypjudge{\oxglobalctx}{\oxemptyctx}{\oxkontctx}{\oxvarctx}
    {\oxabort{\oxdots}}{\oxtref{\oxcprov}{\oxmuta}{\oxtslice{\oxsitype}}}
    {\oxvarctx^\prime}
  }{
    An \oxkey{abort!} expression is well-typed (at any type) via the rule \oxname{T-Abort}.
  }

  %% new type is a subtype of the old type
  \oxpresline{
    \oxtunify[\oxcombine]{\oxemptyctx}{\oxkontctx}{\oxvarctx}{
      \oxtref{\oxcprov}{\oxmuta}{\oxtslice{\oxsitype}}
    }{
      \oxtref{\oxcprov}{\oxmuta}{\oxtslice{\oxsitype}}
    }{\oxvarctx}
  }{
    Immediate by \oxname{RR-Refl}.
  }

  %% new loan context is a sub-context of the old loan context
  \oxpresline{
    \exists \oxvarctx_o. \oxvarctx_i \oxintersect \oxvarctx_o = \oxvarctx_i
  }{
    $\oxvarctx_o = \oxvarctx_i$
  }
\end{preservation}

\oxpreservationcase{T-IndexCopy}{\TIndexCopy}{
  \oxindex{\oxplaceexpr}{\oxexpr_i}
}{\EIndexCopy}{
  \oxvarctx^\prime
}

\begin{preservation}
  %% new store is valid
  \oxpresline{
    \oxstorevalidity{\oxglobalctx}{\oxvarctx^\prime}{\oxstore}
  }{
    Applying \Lemma{lemma:values-dont-change} to the typing derivation (from
    \oxname{T-IndexCopy}) for $\oxexpr$ (which we know is a value from
    \oxname{E-IndexCopy}) gives us
    $\oxsubtypectx{\oxglobalctx}{\oxemptyctx}{\oxvarctx}{\oxvarctx^\prime}$.
    Then, applying \Lemma{lemma:stack-validity-related-envs} to
    $\oxsubtypectx{\oxglobalctx}{\oxemptyctx}{\oxvarctx}{\oxvarctx^\prime}$ and
    $\oxstorevalidity{\oxglobalctx}{\oxvarctx}{\oxstore}$ (from premise) gives us
    $\oxstorevalidity{\oxglobalctx}{\oxvarctx^\prime}{\oxstore}$.
  }

  %% kontctx values are still valid
  \oxpresline{
    \oxwfkontctx{\oxglobalctx}{\oxvarctx^\prime}{\overline{\oxvalue}}{\oxkontctx}
  }{
    Inverting \oxname{WF-Temporaries} gives us $\forall i \in 1 \oxdots n. \;
    \oxtypjudge{\oxglobalctx}{\oxemptyctx}{\oxsitype_1 \oxdotsc \oxsitype_{i-1}}{\oxvarctx}{
      \oxvalue_i
    }{\oxsitype_i}{\oxvarctx}$.

    Applying \Lemma{lemma:values-dont-change} to the typing derivation (from
    \oxname{T-IndexCopy}) for
    $\oxexpr$ (which we know is a value from \oxname{E-IndexCopy}) gives us
    $\oxsubtypectx{\oxglobalctx}{\oxemptyctx}{\oxvarctx}{\oxvarctx^\prime}$.

    Thus, we can apply \Lemma{lemma:value-typing-related-envs} to each of the
    typing judgments to get $\forall i \in 1 \oxdots n. \;
    \oxtypjudge{\oxglobalctx}{\oxemptyctx}{\oxsitype_1 \oxdotsc \oxsitype_{i-1}}{\oxvarctx^\prime}{
      \oxvalue_i
    }{\oxsitype_i}{\oxvarctx^\prime}$.

    Finally, applying \oxname{WF-Temporaries} gives us
    $\oxwfkontctx{\oxglobalctx}{\oxvarctx^\prime}{\overline{\oxvalue}}{\oxkontctx}$.
  }

  %% new expression is well-typed
  \oxpresline{
    \oxtypjudge{\oxglobalctx}{\oxemptyctx}{\oxkontctx}{\oxvarctx^\prime}{
      \oxvalue_{n_i}
    }{\oxsitype}{\oxvarctx^\prime}
  }{
    Applying \Lemma{lemma:values-dont-change} to the typing derivation (from
    \oxname{T-IndexCopy}) for $\oxexpr$ (which we know is a value from
    \oxname{E-IndexCopy}) gives us
    $\oxsubtypectx{\oxglobalctx}{\oxemptyctx}{\oxvarctx}{\oxvarctx^\prime}$.
    Then, applying \Lemma{lemma:stack-validity-related-envs} to
    $\oxsubtypectx{\oxglobalctx}{\oxemptyctx}{\oxvarctx}{\oxvarctx^\prime}$ and
    $\oxstorevalidity{\oxglobalctx}{\oxvarctx}{\oxstore}$ (from premise) gives us
    $\oxstorevalidity{\oxglobalctx}{\oxvarctx^\prime}{\oxstore}$.

    Applying \Lemma{lemma:place-exprs-reduce} to
    $\oxmusafety{\oxemptyctx}{\oxvarctx^\prime}{\oximm}{\oxplaceexpr}{
      \oxset{\overline{\oxloan}}
    }$, $\oxcomputetynoprov{\oxemptyctx}{\oxvarctx^\prime}{\oximm}{\oxplaceexpr}
    {\oxsitype}$, and $\oxstorevalidity{\oxglobalctx}{\oxvarctx^\prime}{\oxstore}$
    % gives us $\oxtypjudge{\oxglobalctx}{\oxemptyctx}{\oxvarctx^\prime}{
    %   \oxarr{\oxvalue_0 \oxdotsc \oxvalue_{\oxnum_i} \oxdotsc \oxvalue_n}
    % }{\oxxitype}{\oxvarctx^\prime}$. Then, by inversion of either \oxname{T-Array} or
    \oxname{T-Slice} (based on whether $\oxxitype = \oxtarr{\oxsitype}{\oxnum}$ or
    $\oxtslice{\oxsitype}$ respectively), we get
    $\forall i \in \oxset{1 \oxdots n}. \;
    \oxtypjudge{\oxglobalctx}{\oxemptyctx}{\oxkontctx}{\oxvarctx^\prime}{
      \oxvalue_i
    }{\oxsitype}{\oxvarctx^\prime}$ (after accounting for the fact that the
    constituent expressions are values and the output environment matches the
    input environment). Thus, we can pick out specifically that
    $\oxtypjudge{\oxglobalctx}{\oxemptyctx}{\oxkontctx}{\oxvarctx^\prime}{
      \oxvalue_i
    }{\oxsitype}{\oxvarctx^\prime}$.
  }

  %% new type is a subtype of the old type
  \oxpresline{
    \oxtunify[\oxcombine]{\oxemptyctx}{\oxkontctx}{\oxvarctx^\prime}{\oxsitype}{\oxsitype}{\oxvarctx^\prime}
  }{
    Immediate by \oxname{RR-Refl}.
  }

  %% new loan context is a sub-context of the old loan context
  \oxpresline{
    \exists \oxvarctx_o. \oxvarctx^\prime \oxintersect \oxvarctx_o = \oxvarctx^\prime
  }{
    $\oxvarctx_o = \oxvarctx^\prime$
  }
\end{preservation}

\oxpreservationcase{T-Seq}{\TSeq}{\oxseq{\oxexpr_1}{\oxexpr_2}}{\ESeq}{
  \oxgcloans{\oxkontctx}{\oxvarctx_1}
}

\begin{preservation}
  %% new store is valid
  \oxpresline{
    \oxstorevalidity{\oxglobalctx}{
      \oxgcloans{\oxkontctx}{\oxvarctx_1}
    }{\oxstore}
  }{
    Applying \Lemma{lemma:values-dont-change} to the typing derivation (from
    \oxname{T-Seq}) for $\oxexpr_1$ (which we know is a value from
    \oxname{E-Seq}) gives us
    $\oxsubtypectx{\oxglobalctx}{\oxemptyctx}{\oxvarctx}{\oxvarctx_1}$.
    Applying \Lemma{lemma:stack-validity-related-envs} with this and
    $\oxstorevalidity{\oxglobalctx}{\oxvarctx}{\oxstore}$ (from
    premise) gives us
    $\oxstorevalidity{\oxglobalctx}{\oxvarctx_1}{\oxstore}$. By definition of
    \oxname{R-Env}, we know that $\oxsubtypectx{\oxglobalctx}{\oxemptyctx}
    {\oxvarctx_1}{\oxgcloans{\oxkontctx}{\oxvarctx_1}}$. Then, we can apply
    \Lemma{lemma:stack-validity-related-envs} to get
    $\oxstorevalidity{\oxglobalctx}{
      \oxgcloans{\oxkontctx}{\oxvarctx_1}
    }{\oxstore}$.
  }

  %% kontctx values are still valid
  \oxpresline{
    \oxwfkontctx{\oxglobalctx}{
      \oxgcloans{\oxkontctx}{\oxvarctx_1}
    }{\overline{\oxvalue}}{\oxkontctx}
  }{
    Inverting \oxname{WF-Temporaries} gives us $\forall i \in 1 \oxdots n. \;
    \oxtypjudge{\oxglobalctx}{\oxemptyctx}{\oxsitype_1 \oxdotsc \oxsitype_{i-1}}{\oxvarctx}{
      \oxvalue_i
    }{\oxsitype_i}{\oxvarctx}$.

    Applying \Lemma{lemma:values-dont-change} to the typing derivation (from
    \oxname{T-Seq}) for
    $\oxexpr_1$ (which we know is a value from \oxname{E-Seq}) gives us
    $\oxsubtypectx{\oxglobalctx}{\oxemptyctx}{\oxvarctx}{\oxvarctx_1}$. Then,
    we note that definitionally $\oxsubtypectx{\oxglobalctx}{\oxemptyctx}{\oxvarctx_1}{
      \oxgcloans{\oxkontctx}{\oxvarctx_1}
    }$ since we can only clear loans that are not present in any of the types.
    Then, by transitivity, we have
    $\oxsubtypectx{\oxglobalctx}{\oxemptyctx}{\oxvarctx}{
      \oxgcloans{\oxkontctx}{\oxvarctx_1}
    }$.

    Thus, we can apply \Lemma{lemma:value-typing-related-envs} to each of the
    typing judgments to get $\forall i \in 1 \oxdots n. \;
    \oxtypjudge{\oxglobalctx}{\oxemptyctx}{\oxsitype_1 \oxdotsc \oxsitype_{i-1}}{
      \oxgcloans{\oxkontctx}{\oxvarctx_1}
    }{
      \oxvalue_i
    }{\oxsitype_i}{
      \oxgcloans{\oxkontctx}{\oxvarctx_1}
    }$.

    Finally, applying \oxname{WF-Temporaries} gives us
    $\oxwfkontctx{\oxglobalctx}{
      \oxgcloans{\oxkontctx}{\oxvarctx_1}
    }{\overline{\oxvalue}}{\oxkontctx}$.
  }

  %% new expression is well-typed
  \oxpresline{
    \oxtypjudge{\oxglobalctx}{\oxemptyctx}{\oxkontctx}{\oxgcloans{\oxkontctx}{\oxvarctx_1}}
    {\oxexpr_2}{\oxsitype_2}{\oxvarctx_2}
  }{
    Immediate from the premise of \oxname{T-Seq}.
  }

  %% new type is a subtype of the old type
  \oxpresline{
    \oxtunify[\oxcombine]{\oxemptyctx}{\oxkontctx}{\oxvarctx_2}
    {\oxsitype_2}{\oxsitype_2}{\oxvarctx_2}
  }{
    Immediate by \oxname{RR-Refl}.
  }

  %% new loan context is a sub-context of the old loan context
  \oxpresline{
    \exists \oxvarctx_o. \oxvarctx_2 \oxintersect \oxvarctx_o = \oxvarctx_2
  }{
    $\oxvarctx_o = \oxvarctx_2$
  }
\end{preservation}

\oxpreservationcasemulti{T-Branch}{\TBranch}{
  \oxbranch{\oxexpr_1}{\oxexpr_2}{\oxexpr_3}
}{\EIfTrue \and \EIfFalse}

\oxpreservationsubcaseheading{E-IfTrue}{
  \oxvarctx_1
}

\begin{preservation}
  %% new store is valid
  \oxpresline{
    \oxstorevalidity{\oxglobalctx}{\oxvarctx_1}{\oxstore}
  }{
    Applying \Lemma{lemma:values-dont-change} to the typing derivation (from
    \oxname{T-Branch}) for $\oxexpr_1$ (which we know is a value from
    \oxname{E-IfTrue}) gives us
    $\oxsubtypectx{\oxglobalctx}{\oxemptyctx}{\oxvarctx}{\oxvarctx_1}$. Then,
    applying \Lemma{lemma:stack-validity-related-envs} with this and
    $\oxstorevalidity{\oxglobalctx}{\oxvarctx}{\oxstore}$ (from
    premise) gives us
    $\oxstorevalidity{\oxglobalctx}{\oxvarctx_1}{\oxstore}$.
  }

  %% kontctx values are still valid
  \oxpresline{
    \oxwfkontctx{\oxglobalctx}{\oxvarctx_1}{\overline{\oxvalue}}{\oxkontctx}
  }{
    Inverting \oxname{WF-Temporaries} gives us $\forall i \in 1 \oxdots n. \;
    \oxtypjudge{\oxglobalctx}{\oxemptyctx}{\oxsitype_1 \oxdotsc \oxsitype_{i-1}}{\oxvarctx}{
      \oxvalue_i
    }{\oxsitype_i}{\oxvarctx}$.

    Applying \Lemma{lemma:values-dont-change} to the typing derivation (from
    \oxname{T-Branch}) for
    $\oxexpr_1$ (which we know is a value from \oxname{E-IfTrue}) gives us
    $\oxsubtypectx{\oxglobalctx}{\oxemptyctx}{\oxvarctx}{\oxvarctx_1}$.

    Thus, we can apply \Lemma{lemma:value-typing-related-envs} to each of the
    typing judgments to get $\forall i \in 1 \oxdots n. \;
    \oxtypjudge{\oxglobalctx}{\oxemptyctx}{\oxsitype_1 \oxdotsc \oxsitype_{i-1}}{\oxvarctx_1}{
      \oxvalue_i
    }{\oxsitype_i}{\oxvarctx_1}$.

    Finally, applying \oxname{WF-Temporaries} gives us
    $\oxwfkontctx{\oxglobalctx}{\oxvarctx_1}{\overline{\oxvalue}}{\oxkontctx}$.
  }

  %% new expression is well-typed
  \oxpresline{
    \oxtypjudge{\oxglobalctx}{\oxemptyctx}{\oxkontctx}{\oxvarctx_1}{\oxexpr_2}
    {\oxsitype}{\oxvarctx_2}
  }{
    Immediate from premise of \oxname{T-Branch}.
  }

  %% new type is a subtype of the old type
  \oxpresline{
    \oxtunify[\oxcombine]{\oxemptyctx}{\oxkontctx}{\oxvarctx_2}
    {\oxsitype_2}{\oxsitype}{\oxvarctx_2^\prime}
  }{
    Immediate from premise of \oxname{T-Branch}.
  }

  %% new loan context is a sub-context of the old loan context
  \oxpresline{
    \exists \oxvarctx_o. \oxvarctx_2^\prime \oxintersect \oxvarctx_o = \oxvarctx^\prime
  }{
    $\oxvarctx_o = \oxvarctx_3^\prime$
  }
\end{preservation}

\oxpreservationsubcaseheading{E-IfFalse}{
  \oxvarctx_1
}

\begin{preservation}
  %% new store is valid
  \oxpresline{
    \oxstorevalidity{\oxglobalctx}{\oxvarctx_1}{\oxstore}
  }{
    Applying \Lemma{lemma:values-dont-change} to the typing derivation (from
    \oxname{T-Branch}) for $\oxexpr_1$ (which we know is a value from
    \oxname{E-IfFalse}) gives us
    $\oxsubtypectx{\oxglobalctx}{\oxemptyctx}{\oxvarctx}{\oxvarctx_1}$. Then,
    applying \Lemma{lemma:stack-validity-related-envs} with this and
    $\oxstorevalidity{\oxglobalctx}{\oxvarctx}{\oxstore}$ (from
    premise) gives us
    $\oxstorevalidity{\oxglobalctx}{\oxvarctx_1}{\oxstore}$.
  }

  %% kontctx values are still valid
  \oxpresline{
    \oxwfkontctx{\oxglobalctx}{\oxvarctx_1}{\overline{\oxvalue}}{\oxkontctx}
  }{
    Inverting \oxname{WF-Temporaries} gives us $\forall i \in 1 \oxdots n. \;
    \oxtypjudge{\oxglobalctx}{\oxemptyctx}{\oxsitype_1 \oxdotsc \oxsitype_{i-1}}{\oxvarctx}{
      \oxvalue_i
    }{\oxsitype_i}{\oxvarctx}$.

    Applying \Lemma{lemma:values-dont-change} to the typing derivation (from
    \oxname{T-Branch}) for
    $\oxexpr_1$ (which we know is a value from \oxname{E-IfFalse}) gives us
    $\oxsubtypectx{\oxglobalctx}{\oxemptyctx}{\oxvarctx}{\oxvarctx_1}$.

    Thus, we can apply \Lemma{lemma:value-typing-related-envs} to each of the
    typing judgments to get $\forall i \in 1 \oxdots n. \;
    \oxtypjudge{\oxglobalctx}{\oxemptyctx}{\oxsitype_1 \oxdotsc \oxsitype_{i-1}}{\oxvarctx_1}{
      \oxvalue_i
    }{\oxsitype_i}{\oxvarctx_1}$.

    Finally, applying \oxname{WF-Temporaries} gives us
    $\oxwfkontctx{\oxglobalctx}{\oxvarctx_1}{\overline{\oxvalue}}{\oxkontctx}$.
  }

  %% new expression is well-typed
  \oxpresline{
    \oxtypjudge{\oxglobalctx}{\oxemptyctx}{\oxkontctx}{\oxvarctx_1}
    {\oxexpr_3}{\oxsitype_3}{\oxvarctx_3}
  }{
    Immediate from premise of \oxname{T-Branch}.
  }

  %% new type is a subtype of the old type
  \oxpresline{
    \oxtunify[\oxcombine]{\oxemptyctx}{\oxkontctx}{\oxvarctx_3}
    {\oxsitype_3}{\oxsitype}{\oxvarctx_3^\prime}
  }{
    Immediate from premise of \oxname{T-Branch}.
  }

  %% new loan context is a sub-context of the old loan context
  \oxpresline{
    \exists \oxvarctx_o. \oxvarctx_3^\prime \oxintersect \oxvarctx_o = \oxvarctx^\prime
  }{
    $\oxvarctx_o = \oxvarctx_2^\prime$ (note that $\oxintersect$ commutes)
  }
\end{preservation}

\oxpreservationcasemulti{T-Match}{\TMatch}{
  \oxmatch{\oxexpr}{\oxid_1}{\oxexpr_1}{\oxid_2}{\oxexpr_2}
}{\EMatchLeft \and \EMatchRight}

\oxpreservationsubcaseheading{E-MatchLeft}{
  \oxextendctx{\oxvarctx^\prime}{
    \oxvarctxentry{\oxid_1}{\oxsitype_l}
  }
}

\begin{preservation}
  %% new store is valid
  \oxpresline{
    \oxstorevalidity{\oxglobalctx}{
      \oxextendctx{\oxvarctx^\prime}{
        \oxvarctxentry{\oxid_1}{\oxsitype_l}
      }
    }{
      \oxextendctx{\oxstore}{
        \oxstoreentry{\oxid_1}{\oxvalue}
      }
    }
  }{
    Applying \Lemma{lemma:values-dont-change} to the typing derivation (from
    \oxname{T-Match}) for $\oxexpr$ (which we know is a value from
    \oxname{E-MatchLeft}) gives us
    $\oxsubtypectx{\oxglobalctx}{\oxemptyctx}{\oxvarctx}{\oxvarctx^\prime}$. Then,
    applying \Lemma{lemma:stack-validity-related-envs} with this and
    $\oxstorevalidity{\oxglobalctx}{\oxvarctx}{\oxstore}$ (from
    premise) gives us
    $\oxstorevalidity{\oxglobalctx}{\oxvarctx^\prime}{\oxstore}$. Then, using
    $\oxtypjudge{\oxglobalctx}{\oxemptyctx}{\oxkontctx}{\oxvarctx}{\oxvalue}{\oxsitype_l}
    {\oxvarctx^\prime}$ (which we get by inversion of \oxname{T-Inl} in the
    derivation for $\oxexpr$) with \oxname{WF-StackFrame} allows us to finally
    conclude $\oxstorevalidity{\oxglobalctx}{
      \oxextendctx{\oxvarctx^\prime}{
        \oxvarctxentry{\oxid_1}{\oxsitype_l}
      }
    }{
      \oxextendctx{\oxstore}{
        \oxstoreentry{\oxid_1}{\oxvalue}
      }
    }$.
  }

  %% kontctx values are still valid
  \oxpresline{
    \oxwfkontctx{\oxglobalctx}{
      \oxextendctx{\oxvarctx^\prime}{
        \oxvarctxentry{\oxid_1}{\oxsitype_l}
      }
    }{\overline{\oxvalue}}{\oxkontctx}
  }{
    Inverting \oxname{WF-Temporaries} gives us $\forall i \in 1 \oxdots n. \;
    \oxtypjudge{\oxglobalctx}{\oxemptyctx}{\oxsitype_1 \oxdotsc \oxsitype_{i-1}}{\oxvarctx}{
      \oxvalue_i
    }{\oxsitype_i}{\oxvarctx}$.

    Applying \Lemma{lemma:values-dont-change} to the typing derivation (from
    \oxname{T-Match}) for
    $\oxexpr$ (which we know is a value from \oxname{E-MatchLeft}) gives us
    $\oxsubtypectx{\oxglobalctx}{\oxemptyctx}{\oxvarctx}{\oxvarctx^\prime}$.

    Thus, we can apply \Lemma{lemma:value-typing-related-envs} to each of the
    typing judgments to get $\forall i \in 1 \oxdots n. \;
    \oxtypjudge{\oxglobalctx}{\oxemptyctx}{\oxsitype_1 \oxdotsc \oxsitype_{i-1}}{\oxvarctx^\prime}{
      \oxvalue_i
    }{\oxsitype_i}{\oxvarctx^\prime}$.

    Applying \Lemma{lemma:values-well-typed-extension} to each of the value
    typing judgments and $\forall \oxcprov \in
    \oxfprovs{\oxtsum{\oxsitype_l}{\oxsitype_r}}. \;
    \oxnotreborrowed{\oxvarctx^\prime}{\oxcprov}$ (from the premise of
    \oxname{T-Match}) gives us
    $\forall i \in 1 \oxdots n. \;
    \oxtypjudge{\oxglobalctx}{\oxemptyctx}{\oxsitype_1 \oxdotsc \oxsitype_{i-1}}{
      \oxextendctx{\oxvarctx^\prime}{
        \oxvarctxentry{\oxid_1}{\oxsitype_l}
      }
    }{
      \oxvalue_i
    }{\oxsitype_i}{
      \oxextendctx{\oxvarctx^\prime}{
        \oxvarctxentry{\oxid_1}{\oxsitype_l}
      }
    }$.

    Finally, applying \oxname{WF-Temporaries} gives us
    $\oxwfkontctx{\oxglobalctx}{
      \oxextendctx{\oxvarctx^\prime}{
        \oxvarctxentry{\oxid_1}{\oxsitype_l}
      }
    }{\overline{\oxvalue}}{\oxkontctx}$.
  }

  %% new expression is well-typed
  \oxpresline{
    \oxtypjudge{\oxglobalctx}{\oxemptyctx}{\oxkontctx}{
      \oxextendctx{\oxvarctx^\prime}{
        \oxvarctxentry{\oxid_1}{\oxsitype_l}
      }
    }{\oxexpr_1}
    {\oxsitype_1}{\oxvarctx_1}
  }{
    Immediate from applying \oxname{T-Shift} to
    $\oxtypjudge{\oxglobalctx}{\oxemptyctx}{\oxkontctx}{
      \oxextendctx{\oxvarctx^\prime}{
        \oxvarctxentry{\oxid_1}{\oxsitype_l}
      }
    }{
      \oxexpr_1
    }{\oxsitype_1}{
      \oxextendctx{\oxvarctx_1}{
        \oxvarctxentry{\oxid_1}{\oxsdtype_l}
      }
    }$ from the premise of \oxname{T-Match}.
  }

  %% new type is a subtype of the old type
  \oxpresline{
    \oxtunify[\oxcombine]{\oxemptyctx}{\oxkontctx}{\oxvarctx_1}
    {\oxsitype_1}{\oxsitype}{\oxvarctx_1^\prime}
  }{
    Immediate from premise of \oxname{T-Match}.
  }

  %% new loan context is a sub-context of the old loan context
  \oxpresline{
    \exists \oxvarctx_o. \oxvarctx_1^\prime \oxintersect \oxvarctx_o = \oxvarctx^\prime
  }{
    $\oxvarctx_o = \oxvarctx_2^\prime$
  }
\end{preservation}

\oxpreservationsubcaseheading{E-MatchRight}{
  \oxextendctx{\oxvarctx^\prime}{
    \oxvarctxentry{\oxid_2}{\oxsitype_r}
  }
}

\begin{preservation}
  %% new store is valid
  \oxpresline{
    \oxstorevalidity{\oxglobalctx}{
      \oxextendctx{\oxvarctx^\prime}{
        \oxvarctxentry{\oxid_2}{\oxsitype_r}
      }
    }{
      \oxextendctx{\oxstore}{
        \oxstoreentry{\oxid_2}{\oxvalue}
      }
    }
  }{
    Applying \Lemma{lemma:values-dont-change} to the typing derivation (from
    \oxname{T-Match}) for $\oxexpr$ (which we know is a value from
    \oxname{E-MatchRight}) gives us
    $\oxsubtypectx{\oxglobalctx}{\oxemptyctx}{\oxvarctx}{\oxvarctx^\prime}$. Then,
    applying \Lemma{lemma:stack-validity-related-envs} with this and
    $\oxstorevalidity{\oxglobalctx}{\oxvarctx}{\oxstore}$ (from
    premise) gives us
    $\oxstorevalidity{\oxglobalctx}{\oxvarctx^\prime}{\oxstore}$. Then, using
    $\oxtypjudge{\oxglobalctx}{\oxemptyctx}{\oxkontctx}{\oxvarctx}{\oxvalue}{\oxsitype_r}
    {\oxvarctx^\prime}$ (which we get by inversion of \oxname{T-Inr} in the
    derivation for $\oxexpr$) with \oxname{WF-StackFrame} allows us to finally
    conclude $\oxstorevalidity{\oxglobalctx}{
      \oxextendctx{\oxvarctx^\prime}{
        \oxvarctxentry{\oxid_2}{\oxsitype_r}
      }
    }{
      \oxextendctx{\oxstore}{
        \oxstoreentry{\oxid_2}{\oxvalue}
      }
    }$.
  }

  %% kontctx values are still valid
  \oxpresline{
    \oxwfkontctx{\oxglobalctx}{
      \oxextendctx{\oxvarctx^\prime}{
        \oxvarctxentry{\oxid_2}{\oxsitype_r}
      }
    }{\overline{\oxvalue}}{\oxkontctx}
  }{
    Inverting \oxname{WF-Temporaries} gives us $\forall i \in 1 \oxdots n. \;
    \oxtypjudge{\oxglobalctx}{\oxemptyctx}{\oxsitype_1 \oxdotsc \oxsitype_{i-1}}{\oxvarctx}{
      \oxvalue_i
    }{\oxsitype_i}{\oxvarctx}$.

    Applying \Lemma{lemma:values-dont-change} to the typing derivation (from
    \oxname{T-Match}) for
    $\oxexpr$ (which we know is a value from \oxname{E-MatchLeft}) gives us
    $\oxsubtypectx{\oxglobalctx}{\oxemptyctx}{\oxvarctx}{\oxvarctx^\prime}$.

    Thus, we can apply \Lemma{lemma:value-typing-related-envs} to each of the
    typing judgments to get $\forall i \in 1 \oxdots n. \;
    \oxtypjudge{\oxglobalctx}{\oxemptyctx}{\oxsitype_1 \oxdotsc \oxsitype_{i-1}}{\oxvarctx^\prime}{
      \oxvalue_i
    }{\oxsitype_i}{\oxvarctx^\prime}$.

    Applying \Lemma{lemma:values-well-typed-extension} to each of the value
    typing judgments and $\forall \oxcprov \in
    \oxfprovs{\oxtsum{\oxsitype_l}{\oxsitype_r}}. \;
    \oxnotreborrowed{\oxvarctx^\prime}{\oxcprov}$ (from the premise of
    \oxname{T-Match}) gives us
    $\forall i \in 1 \oxdots n. \;
    \oxtypjudge{\oxglobalctx}{\oxemptyctx}{\oxsitype_1 \oxdotsc \oxsitype_{i-1}}{
      \oxextendctx{\oxvarctx^\prime}{
        \oxvarctxentry{\oxid_2}{\oxsitype_r}
      }
    }{
      \oxvalue_i
    }{\oxsitype_i}{
      \oxextendctx{\oxvarctx^\prime}{
        \oxvarctxentry{\oxid_2}{\oxsitype_r}
      }
    }$.

    Finally, applying \oxname{WF-Temporaries} gives us
    $\oxwfkontctx{\oxglobalctx}{
      \oxextendctx{\oxvarctx^\prime}{
        \oxvarctxentry{\oxid_2}{\oxsitype_r}
      }
    }{\overline{\oxvalue}}{\oxkontctx}$.
  }

  %% new expression is well-typed
  \oxpresline{
    \oxtypjudge{\oxglobalctx}{\oxemptyctx}{\oxkontctx}{
      \oxextendctx{\oxvarctx^\prime}{
        \oxvarctxentry{\oxid_2}{\oxsitype_r}
      }
    }{\oxexpr_2}
    {\oxsitype_2}{\oxvarctx_2}
  }{
    Immediate from applying \oxname{T-Shift} to
    $\oxtypjudge{\oxglobalctx}{\oxemptyctx}{\oxkontctx}{
      \oxextendctx{\oxvarctx^\prime}{
        \oxvarctxentry{\oxid_2}{\oxsitype_r}
      }
    }{
      \oxexpr_2
    }{\oxsitype_2}{
      \oxextendctx{\oxvarctx_2}{
        \oxvarctxentry{\oxid_2}{\oxsdtype_r}
      }
    }$ from the premise of \oxname{T-Match}.
  }

  %% new type is a subtype of the old type
  \oxpresline{
    \oxtunify[\oxcombine]{\oxemptyctx}{\oxkontctx}{\oxvarctx_2}
    {\oxsitype_2}{\oxsitype}{\oxvarctx_2^\prime}
  }{
    Immediate from premise of \oxname{T-Match}.
  }

  %% new loan context is a sub-context of the old loan context
  \oxpresline{
    \exists \oxvarctx_o. \oxvarctx_2^\prime \oxintersect \oxvarctx_o = \oxvarctx^\prime
  }{
    $\oxvarctx_o = \oxvarctx_1^\prime$ (note that $\oxintersect$ commutes)
  }
\end{preservation}

\oxpreservationcase{T-Assign}{\TAssign}{
  \oxassign{\oxplace}{\oxexpr_a}
}{\EAssign}{
  \oxgcloans{\oxkontctx}{\oxtupdate{\oxvarctx^\prime}{\oxplace}{\oxsitype}}
}

\begin{preservation}
  %% new store is valid
  \oxpresline{
    \oxstorevalidity{\oxglobalctx}{
      \oxgcloans{\oxkontctx}{\oxtupdate{\oxvarctx^\prime}{\oxplace}{\oxsitype}}
    }{
      \oxvupdate{\oxstore}{\oxplace}{\oxvaluectx[\oxvalue]}
    }
  }{
    From our premise, we have
    $\oxstorevalidity{\oxglobalctx}{\oxvarctx}{\oxstore}$. Then, applying
    \Lemma{lemma:values-dont-change} to the typing derivation
    $\oxtypjudge{\oxglobalctx}{\oxemptyctx}{\oxkontctx}{\oxvarctx}{
      \oxvalue
    }{\oxsitype}{\oxvarctx_1}$ (from the premise of \oxname{T-Assign} combined
    with the fact that $\oxexpr$ is a value from \oxname{E-Assign}) gives us
    $\oxsubtypectx{\oxglobalctx}{\oxemptyctx}{\oxvarctx}{\oxvarctx_1}$. Then,
    applying \Lemma{lemma:stack-validity-related-envs} to these two facts gives
    us $\oxstorevalidity{\oxglobalctx}{\oxvarctx_1}{\oxstore}$.

    Finally, we apply \Lemma{lemma:stack-validity-assign} to
    $\oxstorevalidity{\oxglobalctx}{\oxvarctx_1}{\oxstore}$ and
    $\oxlookup{{\oxvarctx_1}}{\oxplace}{\oxsxtype}$ (from premise) and
    $\oxtunify[\oxnoop]{\oxemptyctx}{\oxkontctx}{
      \oxrsub{\oxvarctx_1}{\oxplace}
    }{\oxsitype}{\oxsxtype}{\oxvarctx^\prime}$ (from the premise of
    \oxname{T-Assign}) and
    $\oxmusafety{\oxemptyctx}{\oxvarctx^\prime}{\oxmut}{\oxplace}{
      \oxset{\oxloanpkg{\oxmut}{\oxplace}}
    }$ (from premise of \oxname{T-Assign} and
    $\oxnorm{\oxstore}{\oxplace}{\oxplace}{\oxvaluectx}{\_}$ and $\oxplace =
    \oxproj{\oxid}{\oxpath}$ (both from the premise of \oxname{E-Assign}),
    $\oxtypjudge{\oxglobalctx}{\oxemptyctx}{\oxkontctx}{
      \oxrsub{\oxvarctx_1}{\oxplaceexpr} }{\oxvalue}{\oxsitype}{
      \oxrsub{\oxvarctx_1}{\oxplaceexpr} }$ (from above) and
    $\oxstorevalidity{\oxglobalctx}{
      \oxgcloans{\oxkontctx}{\oxtupdate{\oxvarctx^\prime}{\oxplace}{\oxsitype}}
    }{
      \oxvupdate{\oxstore}{\oxid}{\oxvaluectx[\oxvalue]}
    }$.
  }

  %% kontctx values are still valid
  \oxpresline{
    \oxwfkontctx{\oxglobalctx}{
      \oxgcloans{\oxkontctx}{\oxtupdate{\oxvarctx^\prime}{\oxplace}{\oxsitype}}
    }{\overline{\oxvalue}}{\oxkontctx}
  }{
    Inverting \oxname{WF-Temporaries} gives us $\forall i \in 1 \oxdots n. \;
    \oxtypjudge{\oxglobalctx}{\oxemptyctx}{\oxsitype_1 \oxdotsc \oxsitype_{i-1}}{\oxvarctx}{
      \oxvalue_i
    }{\oxsitype_i}{\oxvarctx}$.

    Applying \Lemma{lemma:values-dont-change} to the typing derivation (from
    \oxname{T-Let}) for
    $\oxexpr_1$ (which we know is a value from \oxname{E-Let}) gives us
    $\oxsubtypectx{\oxglobalctx}{\oxemptyctx}{\oxvarctx}{\oxvarctx_1}$.
    Then, we can apply \Lemma{lemma:value-typing-related-envs} to each of the
    typing judgments to get $\forall i \in 1 \oxdots n. \;
    \oxtypjudge{\oxglobalctx}{\oxemptyctx}{\oxsitype_1 \oxdotsc \oxsitype_{i-1}}{\oxvarctx_1}{
      \oxvalue_i
    }{\oxsitype_i}{\oxvarctx_1}$.

    Then, we apply \Lemma{lemma:values-well-typed-assignment} to each of these
    typing judgments to get $\forall i \in 1 \oxdots n. \;
    \oxtypjudge{\oxglobalctx}{\oxemptyctx}{\oxsitype_1 \oxdotsc \oxsitype_{i-1}}{
      \oxgcloans{\oxkontctx}{\oxtupdate{\oxvarctx^\prime}{\oxplace}{\oxsitype}}
    }{
      \oxvalue_i
    }{\oxsitype_i}{
      \oxgcloans{\oxkontctx}{\oxtupdate{\oxvarctx^\prime}{\oxplace}{\oxsitype}}
    }$. After which we can apply \oxname{WF-Temporaries} to conclude
    $\oxwfkontctx{\oxglobalctx}{
      \oxgcloans{\oxkontctx}{\oxtupdate{\oxvarctx^\prime}{\oxplace}{\oxsitype}}
    }{\overline{\oxvalue}}{\oxkontctx}$.
  }

  %% new expression is well-typed
  \oxpresline{
    \oxtypjudge{\oxglobalctx}{\oxemptyctx}{\oxkontctx}{
      \oxvarctx_i
    }{\oxunit}{\oxtunit}{
      \oxvarctx_i
    }
  }{
    Immediate by \oxname{T-Unit}.
  }

  %% new type is a subtype of the old type
  \oxpresline{
    \oxtunify[\oxcombine]{\oxemptyctx}{\oxkontctx}{
      \oxvarctx_i
    }{\oxtunit}{\oxtunit}{
      \oxvarctx_i
    }
  }{
    Immediate by \oxname{RR-Refl}.
  }

  %% new loan context is a sub-context of the old loan context
  \oxpresline{
    \exists \oxvarctx_o.
    \oxvarctx_i
    \oxintersect \oxvarctx_o =
    \oxvarctx_i
  }{
    $\oxvarctx_o = \oxvarctx_i =
    \oxgcloans{\oxkontctx}{\oxtupdate{\oxvarctx^\prime}{\oxplace}{\oxsitype}}$
  }
\end{preservation}

\oxpreservationcase{T-AssignDeref}{\TAssignDeref}{
  \oxassign{\oxplaceexpr}{\oxexpr_a}
}{\EAssign}{
  \oxvarctx^\prime
}

\begin{preservation}
  %% new store is valid
  \oxpresline{
    \oxstorevalidity{\oxglobalctx}{\oxvarctx^\prime}{
      \oxvupdate{\oxstore}{\oxplace}{\oxvaluectx[\oxvalue]}
    }
  }{
    From our premise, we have
    $\oxstorevalidity{\oxglobalctx}{\oxvarctx}{\oxstore}$. Then, applying
    \Lemma{lemma:values-dont-change} to the typing derivation
    $\oxtypjudge{\oxglobalctx}{\oxemptyctx}{\oxkontctx}{\oxvarctx}{
      \oxvalue
    }{\oxsitype}{\oxvarctx_1}$ (from the premise of \oxname{T-AssignDeref} combined
    with the fact that $\oxexpr$ is a value from \oxname{E-Assign}) gives us
    $\oxsubtypectx{\oxglobalctx}{\oxemptyctx}{\oxvarctx}{\oxvarctx_1}$. Then,
    applying \Lemma{lemma:stack-validity-related-envs} to these two facts gives
    us $\oxstorevalidity{\oxglobalctx}{\oxvarctx_1}{\oxstore}$.

    Then, applying \Lemma{lemma:stack-validity-region-rewriting} to
    $\oxstorevalidity{\oxglobalctx}{\oxvarctx_1}{\oxstore}$ and
    $\oxtunify[\oxcombine]{\oxemptyctx}{\oxkontctx}{\oxvarctx_1}
    {\oxsitype_n}{\oxsitype_o}{\oxvarctx^\prime}$ (from the premise of
    \oxname{T-AssignDeref}) gives us
    $\oxstorevalidity{\oxglobalctx}{\oxvarctx^\prime}{\oxstore}$.
  }

  %% kontctx values are still valid
  \oxpresline{
    \oxwfkontctx{\oxglobalctx}{
      \oxvarctx^\prime
    }{\overline{\oxvalue}}{\oxkontctx}
  }{
    Inverting \oxname{WF-Temporaries} gives us $\forall i \in 1 \oxdots n. \;
    \oxtypjudge{\oxglobalctx}{\oxemptyctx}{\oxsitype_1 \oxdotsc \oxsitype_{i-1}}{\oxvarctx}{
      \oxvalue_i
    }{\oxsitype_i}{\oxvarctx}$.

    Applying \Lemma{lemma:values-dont-change} to the typing derivation (from
    \oxname{T-AssignDeref}) for
    $\oxexpr$ (which we know is a value from \oxname{E-Assign}) gives us
    $\oxsubtypectx{\oxglobalctx}{\oxemptyctx}{\oxvarctx}{\oxvarctx_1}$.
    Then, we can apply \Lemma{lemma:value-typing-related-envs} to each of the
    typing judgments to get $\forall i \in 1 \oxdots n. \;
    \oxtypjudge{\oxglobalctx}{\oxemptyctx}{\oxsitype_1 \oxdotsc \oxsitype_{i-1}}{\oxvarctx_1}{
      \oxvalue_i
    }{\oxsitype_i}{\oxvarctx_1}$.

    Then, applying \Lemma{lemma:subtyping-value-typing} to each of these
    derivations along with
    $\oxtunify[\oxcombine]{\oxemptyctx}{\oxkontctx}{\oxvarctx_1}
    {\oxsitype_n}{\oxsitype_o}{\oxvarctx^\prime}$ (from the premise of
    \oxname{T-AssignDeref}) gives us $\forall i \in 1 \oxdots n. \;
    \oxtypjudge{\oxglobalctx}{\oxemptyctx}{\oxsitype_1 \oxdotsc \oxsitype_{i-1}}{
      \oxvarctx^\prime
    }{\oxvalue_i}{\oxsitype_i}{\oxvarctx^\prime}$. We can then apply
    \oxname{WF-Temporaries} to conclude
    $\oxwfkontctx{\oxglobalctx}{
      \oxvarctx^\prime
    }{\overline{\oxvalue}}{\oxkontctx}$.
  }

  %% new expression is well-typed
  \oxpresline{
    \oxtypjudge{\oxglobalctx}{\oxemptyctx}{\oxkontctx}{
      \oxvarctx^\prime
    }{\oxunit}{\oxtunit}{
      \oxvarctx^\prime
    }
  }{
    Immediate by \oxname{T-Unit}.
  }

  %% new type is a subtype of the old type
  \oxpresline{
    \oxtunify[\oxcombine]{\oxemptyctx}{\oxkontctx}{
      \oxvarctx^\prime
    }{\oxtunit}{\oxtunit}{
      \oxvarctx^\prime
    }
  }{
    Immediate by \oxname{RR-Refl}.
  }

  %% new loan context is a sub-context of the old loan context
  \oxpresline{
    \exists \oxvarctx_o. \oxvarctx^\prime \oxintersect
    \oxvarctx_o = \oxvarctx^\prime
  }{
    $\oxvarctx_o = \oxvarctx^\prime$
  }
\end{preservation}

\oxpreservationcase{T-Let}{\TLet}{
  \oxlet{\oxid}{\oxsitype}{\oxexpr_1}{\oxexpr_2}
}{\ELet}{
  \oxgcloans{\oxkontctx}{
    \oxextendctx{\oxvarctx_1^\prime}{
      \oxvarctxentry{\oxid}{\oxsitype_a}
    }
  }
}

\begin{preservation}
  %% new store is valid
  \oxpresline{
    \oxstorevalidity{\oxglobalctx}{
      \oxgcloans{\oxkontctx}{
        \oxextendctx{\oxvarctx_1^\prime}{
          \oxvarctxentry{\oxid}{\oxsitype_a}
        }
      }
    }{
      \oxextendctx{\oxstore}{\oxstoreentry{\oxid}{\oxvalue}}
    }
  }{
    Applying \Lemma{lemma:values-dont-change} to the typing derivation (from
    \oxname{T-Let}) for $\oxexpr_1$ (which we know is a value from
    \oxname{E-Let}) gives us
    $\oxsubtypectx{\oxglobalctx}{\oxemptyctx}{\oxvarctx}{\oxvarctx_1}$.
    Then, we can apply \Lemma{lemma:stack-validity-related-envs} to get
    $\oxstorevalidity{\oxglobalctx}{\oxvarctx_1}{\oxstore}$.
    Then,
    applying \Lemma{lemma:stack-validity-region-rewriting} to
    $\oxtunify[\oxcombine]{\oxemptyctx}{\oxkontctx}{\oxvarctx_1}{\oxsitype_1}{\oxsitype_a}
    {\oxvarctx_1^\prime}$ (from premise of \oxname{T-Let}) gives us
    $\oxstorevalidity{\oxglobalctx}{\oxvarctx_1^\prime}{\oxstore}$. We can also
    apply \Lemma{lemma:values-typed-supertype} to $\oxtypjudge{\oxglobalctx}
    {\oxemptyctx}{\oxkontctx}{\oxvarctx}{\oxvalue}{\oxsitype_1}{\oxvarctx_1}$ (from premise
    of \oxname{T-Let}) and $\oxtunify[\oxcombine]{\oxemptyctx}{\oxkontctx}{\oxvarctx_1}
    {\oxsitype_1}{\oxsitype_a}{\oxvarctx_1^\prime}$ (from premise of
    \oxname{T-Let}) gives us $\oxtypjudge{\oxglobalctx}{\oxemptyctx}{\oxkontctx}
    {\oxvarctx^\prime}{\oxexpr_1}{\oxsitype_a}{\oxvarctx^\prime}$

    Then, apply \Lemma{lemma:store-validity-extension} to
    $\oxstorevalidity{\oxglobalctx}{\oxvarctx_1^\prime}{\oxstore}$ and
    $\oxtypjudge{\oxglobalctx}{\oxemptyctx}{\oxkontctx}{\oxvarctx^\prime}
    {\oxexpr_1}{\oxsitype_a}{\oxvarctx^\prime}$ gives us
    $\oxstorevalidity{\oxglobalctx}{
      \oxextendctx{\oxvarctx_1^\prime}{
        \oxvarctxentry{\oxid}{\oxsitype_a}
      }
    }{
      \oxextendctx{\oxstore}{\oxstoreentry{\oxid}{\oxvalue}}
    }$. By definition of
    \oxname{R-Env}, we know that $\oxsubtypectx{\oxglobalctx}{\oxemptyctx}{
      \oxextendctx{\oxvarctx_1^\prime}{
        \oxvarctxentry{\oxid}{\oxsitype_a}
      }
    }{
      \oxgcloans{\oxkontctx}{
        \oxextendctx{\oxvarctx_1^\prime}{
          \oxvarctxentry{\oxid}{\oxsitype_a}
        }
      }
    }$. Then, we can apply
    \Lemma{lemma:stack-validity-related-envs} to conclude
    $\oxstorevalidity{\oxglobalctx}{
      \oxgcloans{\oxkontctx}{
        \oxextendctx{\oxvarctx_1^\prime}{
          \oxvarctxentry{\oxid}{\oxsitype_a}
        }
      }
    }{
      \oxextendctx{\oxstore}{\oxstoreentry{\oxid}{\oxvalue}}
    }$.
  }

  %% kontctx values are still valid
  \oxpresline{
    \oxwfkontctx{\oxglobalctx}{
      \oxgcloans{\oxkontctx}{
        \oxextendctx{\oxvarctx_1^\prime}{
          \oxvarctxentry{\oxid}{\oxsitype_a}
        }
      }
    }{\overline{\oxvalue}}{\oxkontctx}
  }{
    Inverting \oxname{WF-Temporaries} gives us $\forall i \in 1 \oxdots n. \;
    \oxtypjudge{\oxglobalctx}{\oxemptyctx}{\oxsitype_1 \oxdotsc \oxsitype_{i-1}}{\oxvarctx}{
      \oxvalue_i
    }{\oxsitype_i}{\oxvarctx}$.

    Applying \Lemma{lemma:values-dont-change} to the typing derivation (from
    \oxname{T-Let}) for
    $\oxexpr_1$ (which we know is a value from \oxname{E-Let}) gives us
    $\oxsubtypectx{\oxglobalctx}{\oxemptyctx}{\oxvarctx}{\oxvarctx_1}$.
    Then, we can apply \Lemma{lemma:value-typing-related-envs} to each of the
    typing judgments to get $\forall i \in 1 \oxdots n. \;
    \oxtypjudge{\oxglobalctx}{\oxemptyctx}{\oxsitype_1 \oxdotsc \oxsitype_{i-1}}{\oxvarctx_1}{
      \oxvalue_i
    }{\oxsitype_i}{\oxvarctx_1}$.

    Applying \Lemma{lemma:values-well-typed-extension} to each of the value
    typing judgments and $\forall \oxcprov \in \oxfprovs{\oxsitype_a}. \;
    \oxnotreborrowed{\oxvarctx_1^\prime}{\oxcprov}$ (from the premise of
    \oxname{T-Let}) gives us
    $\forall i \in 1 \oxdots n. \;
    \oxtypjudge{\oxglobalctx}{\oxemptyctx}{\oxsitype_1 \oxdotsc \oxsitype_{i-1}}{
      \oxextendctx{\oxvarctx_1}{
        \oxvarctxentry{\oxid}{\oxsitype_a}
      }
    }{
      \oxvalue_i
    }{\oxsitype_i}{
      \oxextendctx{\oxvarctx_1}{
        \oxvarctxentry{\oxid}{\oxsitype_a}
      }
    }$.

    Then, we note that definitionally
    $\oxsubtypectx{\oxglobalctx}{\oxemptyctx}{
      \oxextendctx{\oxvarctx_1}{
        \oxvarctxentry{\oxid}{\oxsitype_a}
      }
    }{
      \oxgcloans{\oxkontctx}{
        \oxextendctx{\oxvarctx_1}{
          \oxvarctxentry{\oxid}{\oxsitype_a}
        }
      }
    }$ since we can only clear loans that are not present in any of the types.
    We can use this with each typing derivation in \Lemma{lemma:values-dont-change},
    and then finally apply \oxname{WF-Temporaries} to get
    $\oxwfkontctx{\oxglobalctx}{
      \oxgcloans{\oxkontctx}{
        \oxextendctx{\oxvarctx_1}{
          \oxvarctxentry{\oxid}{\oxsitype_a}
        }
      }
    }{\overline{\oxvalue}}{\oxkontctx}$.
  }

  %% new expression is well-typed
  \oxpresline{
    \oxtypjudge{\oxglobalctx}{\oxemptyctx}{\oxkontctx}{
      \oxgcloans{\oxkontctx}{
        \oxextendctx{\oxvarctx_1^\prime}{
          \oxvarctxentry{\oxid}{\oxsitype_a}
        }
      }
    }{\oxshift{\oxexpr}}{\oxsitype_2}{\oxvarctx_2}
  }{
    Immediate by applying \oxname{T-Shift} to the derivation
    \oxtypjudge{\oxglobalctx}{\oxemptyctx}{\oxkontctx}{
      \oxgcloans{\oxkontctx}{
        \oxextendctx{\oxvarctx_1^\prime}{
          \oxvarctxentry{\oxid}{\oxsitype_a}
        }
      }
    }{\oxexpr}{\oxsitype_2}{
      \oxextendctx{\oxvarctx_2}{
        \oxvarctxentry{\oxid}{\oxsdtype}
      }
    } (from premise of \oxname{T-Let}).
  }

  %% new type is a subtype of the old type
  \oxpresline{
    \oxtunify[\oxcombine]{\oxemptyctx}{\oxkontctx}{\oxvarctx_2}
    {\oxsitype_2}{\oxsitype_2}{\oxvarctx_2}
  }{
    Immediate by \oxname{RR-Refl}.
  }

  %% new loan context is a sub-context of the old loan context
  \oxpresline{
    \exists \oxvarctx_o. \oxvarctx_2 \oxintersect \oxvarctx_o = \oxvarctx_2
  }{
    $\oxvarctx_o = \oxvarctx_2$
  }
\end{preservation}

\oxpreservationcase{T-LetRegion}{\TLetProv}{
  \oxletrgn{\oxcprov}{\oxexpr}
}{\ELetProv}{
  \oxvarctx^\prime
}

\begin{preservation}
  %% new store is valid
  \oxpresline{
    \oxstorevalidity{\oxglobalctx}{\oxvarctx^\prime}{\oxstore}
  }{
    Applying \Lemma{lemma:values-dont-change} to the typing derivation (from
    \oxname{T-LetRegion}) for $\oxexpr$ (which we know is a value from
    \oxname{E-LetRegion}) gives us
    $\oxsubtypectx{\oxglobalctx}{\oxemptyctx}{\oxvarctx}{\oxvarctx^\prime}$. Then,
    applying \Lemma{lemma:stack-validity-related-envs} with this and
    $\oxstorevalidity{\oxglobalctx}{\oxvarctx}{\oxstore}$ (from
    premise) gives us
    $\oxstorevalidity{\oxglobalctx}{\oxvarctx^\prime}{\oxstore}$.
  }

  %% kontctx values are still valid
  \oxpresline{
    \oxwfkontctx{\oxglobalctx}{\oxvarctx^\prime}{\overline{\oxvalue}}{\oxkontctx}
  }{
    Inverting \oxname{WF-Temporaries} gives us $\forall i \in 1 \oxdots n. \;
    \oxtypjudge{\oxglobalctx}{\oxemptyctx}{\oxsitype_1 \oxdotsc \oxsitype_{i-1}}{\oxvarctx}{
      \oxvalue_i
    }{\oxsitype_i}{\oxvarctx}$.

    Applying \Lemma{lemma:values-dont-change} to the typing derivation (from
    \oxname{T-LetRegion}) for
    $\oxexpr$ (which we know is a value from \oxname{E-LetRegion}) gives us
    $\oxsubtypectx{\oxglobalctx}{\oxemptyctx}{\oxvarctx}{\oxvarctx^\prime}$.

    Thus, we can apply \Lemma{lemma:value-typing-related-envs} to each of the
    typing judgments to get $\forall i \in 1 \oxdots n. \;
    \oxtypjudge{\oxglobalctx}{\oxemptyctx}{\oxsitype_1 \oxdotsc \oxsitype_{i-1}}{\oxvarctx^\prime}{
      \oxvalue_i
    }{\oxsitype_i}{\oxvarctx^\prime}$.

    Finally, applying \oxname{WF-Temporaries} gives us
    $\oxwfkontctx{\oxglobalctx}{\oxvarctx^\prime}{\overline{\oxvalue}}{\oxkontctx}$.
  }

  %% new expression is well-typed
  \oxpresline{
    \oxtypjudge{\oxglobalctx}{\oxemptyctx}{\oxkontctx}{\oxvarctx^\prime}{\oxvalue}
    {\oxsitype}{\oxvarctx^\prime}
  }{
    We know from \oxname{E-LetRegion} that $\oxexpr$ is a value $\oxvalue$. Thus,
    we can apply \Lemma{lemma:value-typing-output} to
    $\oxtypjudge{\oxglobalctx}{\oxemptyctx}{\oxkontctx}{\oxvarctx}{\oxvalue}{\oxsitype}{
      \oxextendctx{\oxvarctx^\prime}{
        \oxloanctxentry{\oxcprov}{\oxset{\overline{\oxloan}}}
      }
    }$ to get $\oxtypjudge{\oxglobalctx}{\oxemptyctx}{\oxkontctx}{
      \oxextendctx{\oxvarctx^\prime}{
        \oxloanctxentry{\oxcprov}{\oxset{\overline{\oxloan}}}
      }
    }{\oxvalue}{\oxsitype}{
      \oxextendctx{\oxvarctx^\prime}{
        \oxloanctxentry{\oxcprov}{\oxset{\overline{\oxloan}}}
      }
    }$.

    We now wish to show that
    $\oxtypjudge{\oxglobalctx}{\oxemptyctx}{\oxkontctx}{\oxvarctx^\prime}{\oxvalue}
    {\oxsitype}{\oxvarctx^\prime}$. By inspecting the grammar of values and
    their typing rules, we know that the only values who depend on the
    context are pointers and closure values. But by inversion on
    $\oxtypjudge{\oxglobalctx}{\oxemptyctx}{\oxkontctx}{\oxvarctx}{
      \oxletrgn{\oxcprov}{\oxvalue}
    }{\oxsitype}{\oxvarctx^\prime}$, we know that
    $\oxtypevalidity{\oxglobalctx}{\oxemptyctx}{\oxvarctx^\prime}{\oxsitype}$.
    Since the type is valid without the frame, we know that the values
    cannot depend on that frame. Thus, we can conclude
    $\oxtypjudge{\oxglobalctx}{\oxemptyctx}{\oxkontctx}{\oxvarctx^\prime}{\oxvalue}
    {\oxsitype}{\oxvarctx^\prime}$.
  }

  %% new type is a subtype of the old type
  \oxpresline{
    \oxtunify[\oxcombine]{\oxemptyctx}{\oxkontctx}{\oxvarctx^\prime}
    {\oxsitype}{\oxsitype}{\oxvarctx^\prime}
  }{
    Immediate by \oxname{RR-Refl}.
  }

  %% new loan context is a sub-context of the old loan context
  \oxpresline{
    \exists \oxvarctx_o. \oxvarctx^\prime \oxintersect \oxvarctx_o = \oxvarctx^\prime
  }{
    $\oxvarctx_o = \oxvarctx^\prime$
  }
\end{preservation}

\oxpreservationcase{T-While}{\TWhile}{
  \oxwhile{\oxexpr_1}{\oxexpr_2}
}{\EWhile}{
  \oxvarctx
}

\begin{preservation}
  %% new store is valid
  \oxpresline{
    \oxstorevalidity{\oxglobalctx}{\oxvarctx}{\oxstore}
  }{
    Immediate from our premise.
  }

  %% kontctx values are still valid
  \oxpresline{
    \oxwfkontctx{\oxglobalctx}{\oxvarctx}{\overline{\oxvalue}}{\oxkontctx}
  }{
    Immediate from our premise.
  }

  %% new expression is well-typed
  \oxpresline{
    \oxtypjudge{\oxglobalctx}{\oxemptyctx}{\oxkontctx}{\oxvarctx}{
      \oxexpr^\prime
      % \oxbranch{\oxexpr_1}{
      %   \oxseq{\oxexpr_2}{
      %     \oxwhile{\oxexpr_1}{\oxexpr_2}
      %   }
      % }{\oxunit}
    }{\oxtunit}{\oxvarctx_2}
  }{
    We would like to build a derivation to show that the expression
    $\oxbranch{\oxexpr_1}{
      \oxseq{\oxexpr_2}{
        \oxwhile{\oxexpr_1}{\oxexpr_2}
      }
    }{\oxunit}$ is well-typed. We thus start by applying \oxname{T-Branch}.

    This requires us to show three things. First, $\oxtypjudge{\oxglobalctx}
    {\oxemptyctx}{\oxkontctx}{\oxvarctx}{\oxexpr_1}{\oxtbool}{\oxvarctx_1}$ which we have
    from the premise of \oxname{T-While}. Second, $\oxtypjudge{\oxglobalctx}
    {\oxemptyctx}{\oxkontctx}{\oxvarctx_1}{
      \oxseq{\oxexpr_2}{
        \oxwhile{\oxexpr_1}{\oxexpr_2}
      }
    }{\oxtunit}{\oxvarctx_2}$. We build this by applying \oxname{T-Seq} to
    $\oxtypjudge{\oxglobalctx}{\oxemptyctx}{\oxkontctx}{\oxvarctx_1}{\oxexpr_2}
    {\oxtunit}{\oxvarctx_2}$ and $\oxtypjudge{\oxglobalctx}{\oxemptyctx}{\oxkontctx}
    {\oxvarctx_2}{\oxwhile{\oxexpr_1}{\oxexpr_2}}{\oxvarctx_2}$. The former is
    directly in the premise of \oxname{T-While} and the latter can be built by
    applying \oxname{T-While} to
    $\oxtypjudge{\oxglobalctx}{\oxemptyctx}{\oxkontctx}{\oxvarctx_2}{
      \oxexpr_1
    }{\oxtbool}{\oxvarctx_2}$ and
    $\oxtypjudge{\oxglobalctx}{\oxemptyctx}{\oxkontctx}{\oxvarctx_2}{
      \oxexpr_2
    }{\oxtunit}{\oxvarctx_2}$, both from the premise of our original
    \oxname{T-While}. Finally, we need to show
    $\oxtypjudge{\oxglobalctx}{\oxemptyctx}{\oxkontctx}{\oxvarctx_2}{
      \oxunit
    }{\oxtunit}{\oxvarctx_2}$, which is immediate from \oxname{T-Unit}.
  }

  %% new type is a subtype of the old type
  \oxpresline{
    \oxtunify[\oxcombine]{\oxemptyctx}{\oxkontctx}{\oxvarctx_2}
    {\oxtunit}{\oxtunit}{\oxvarctx_2}
  }{
    Immediate by \oxname{RR-Refl}.
  }

  %% new loan context is a sub-context of the old loan context
  \oxpresline{
    \exists \oxvarctx_o. \oxvarctx_2 \oxintersect \oxvarctx_o = \oxvarctx_2
  }{
    $\oxvarctx_o = \oxvarctx_2$
  }
\end{preservation}

\oxpreservationcasemulti{T-ForArray}{\TForArray}{
  \oxfor{\oxid}{\oxexpr_1}{\oxexpr_2}
}{
  \EForArray \and \EForEmptyArray
}

\oxpreservationsubcaseheading{E-ForArray}{
  \oxextendctx{\oxvarctx_1}{\oxvarctxentry{\oxid}{\oxsitype}}
}

\begin{preservation}
  %% new store is valid
  \oxpresline{
    \oxstorevalidity{\oxglobalctx}{
      \oxextendctx{\oxvarctx_1}{\oxvarctxentry{\oxid}{\oxsitype}}
    }{
      \oxextendctx{\oxstore}{
        \oxstoreentry{\oxid}{\oxvalue_0}
      }
    }
  }{
    Applying \Lemma{lemma:values-dont-change} to the typing derivation (from
    \oxname{T-ForArray}) for $\oxexpr_1$ (which we know is a value from
    \oxname{E-ForArray}) gives us
    $\oxsubtypectx{\oxglobalctx}{\oxemptyctx}{\oxvarctx}{\oxvarctx_1}$. Then,
    we can apply \Lemma{lemma:stack-validity-related-envs} to get
    $\oxstorevalidity{\oxglobalctx}{\oxvarctx_1}{\oxstore}$. Applying
    \Lemma{lemma:value-typing-output} to the derivation
    $\oxtypjudge{\oxglobalctx}{\oxemptyctx}{\oxkontctx}{\oxvarctx}
    {\oxarr{\oxvalue_0 \oxdotsc \oxvalue_n}}{\oxtarr{\oxsitype}{\oxnum}}{\oxvarctx_1}$
    gives us $\oxtypjudge{\oxglobalctx}{\oxemptyctx}{\oxkontctx}{\oxvarctx_1}
    {\oxarr{\oxvalue_0 \oxdotsc \oxvalue_n}}{\oxtarr{\oxsitype}{\oxnum}}{\oxvarctx_1}$.
    Then, using inversion (of \oxname{T-Array}), we get
    $\oxtypjudge{\oxglobalctx}{\oxemptyctx}{\oxkontctx}{\oxvarctx_1}
    {\oxvalue_0}{\oxsitype}{\oxvarctx_1}$. Finally, applying
    \Lemma{lemma:store-validity-extension} to
    $\oxstorevalidity{\oxglobalctx}{\oxvarctx_1}{\oxstore}$ and
    $\oxtypjudge{\oxglobalctx}{\oxemptyctx}{\oxkontctx}{\oxvarctx_1}
    {\oxvalue_0}{\oxsitype}{\oxvarctx_1}$ gives us
    $\oxstorevalidity{\oxglobalctx}{
      \oxextendctx{\oxvarctx_1}{
        \oxvarctxentry{\oxid}{\oxsitype}
      }
    }{
      \oxextendctx{\oxstore}{\oxstoreentry{\oxid}{\oxvalue_0}}
    }$.
  }

  %% kontctx values are still valid
  \oxpresline{
    \oxwfkontctx{\oxglobalctx}{
      \oxextendctx{\oxvarctx_1}{
        \oxvarctxentry{\oxid}{\oxsitype}
      }
    }{\overline{\oxvalue}}{\oxkontctx}
  }{
    Inverting \oxname{WF-Temporaries} gives us $\forall i \in 1 \oxdots n. \;
    \oxtypjudge{\oxglobalctx}{\oxemptyctx}{\oxsitype_1 \oxdotsc \oxsitype_{i-1}}{\oxvarctx}{
      \oxvalue_i
    }{\oxsitype_i}{\oxvarctx}$.

    Applying \Lemma{lemma:values-dont-change} to the typing derivation (from
    \oxname{T-ForArray}) for
    $\oxexpr_1$ (which we know is a value from \oxname{E-ForArray}) gives us
    $\oxsubtypectx{\oxglobalctx}{\oxemptyctx}{\oxvarctx}{\oxvarctx_1}$.

    Thus, we can apply \Lemma{lemma:value-typing-related-envs} to each of the
    typing judgments to get $\forall i \in 1 \oxdots n. \;
    \oxtypjudge{\oxglobalctx}{\oxemptyctx}{\oxemptyctx}{\oxvarctx_1}{
      \oxvalue_i
    }{\oxsitype_i}{\oxvarctx_1}$.

    Applying \Lemma{lemma:values-well-typed-extension} to each of the value
    typing judgments and $\forall \oxcprov \in \oxfprovs{\oxsitype}. \;
    \oxnotreborrowed{\oxvarctx_1}{\oxcprov}$ (from the premise of
    \oxname{T-ForArray}) gives us
    $\forall i \in 1 \oxdots n. \;
    \oxtypjudge{\oxglobalctx}{\oxemptyctx}{\oxemptyctx}{
      \oxextendctx{\oxvarctx_1}{
        \oxvarctxentry{\oxid}{\oxsitype}
      }
    }{
      \oxvalue_i
    }{\oxsitype_i}{
      \oxextendctx{\oxvarctx_1}{
        \oxvarctxentry{\oxid}{\oxsitype}
      }
    }$.
  }

  %% new expression is well-typed
  \oxpresline{
    \oxtypjudge{\oxglobalctx}{\oxemptyctx}{\oxkontctx}{
      \oxextendctx{\oxvarctx_1}{\oxvarctxentry{\oxid}{\oxsitype}}
    }{\oxexpr^\prime}{\oxtunit}{\oxvarctx_1}
  }{
    We need to build a derivation for the expression $\oxseq{\oxshift{\oxexpr}}{
      \oxfor{\oxid}{\oxarr{\oxvalue_1 \oxdotsc \oxvalue_n}}{\oxexpr}
    }$. The bottom of this derivation will be \oxname{T-Seq} which requires us
    to show that $\oxtypjudge{\oxglobalctx}{\oxemptyctx}{\oxkontctx}{
      \oxextendctx{\oxvarctx_1}{
        \oxvarctxentry{\oxid}{\oxsitype}
      }
    }{\oxshift{\oxexpr}}{\oxtunit}{\oxvarctx_1}$ and that
    $\oxtypjudge{\oxglobalctx}{\oxemptyctx}{\oxkontctx}{\oxvarctx_1}{
      \oxfor{\oxid}{\oxarr{\oxvalue_1 \oxdotsc \oxvalue_n}}{\oxexpr}
    }{\oxtunit}{\oxvarctx_1}$.

    To show $\oxtypjudge{\oxglobalctx}{\oxemptyctx}{\oxkontctx}{
      \oxextendctx{\oxvarctx_1}{
        \oxvarctxentry{\oxid}{\oxsitype}
      }
    }{\oxshift{\oxexpr}}{\oxtunit}{\oxvarctx_1}$, we apply \oxname{T-Shift} to
    $\oxtypjudge{\oxglobalctx}{\oxemptyctx}{\oxkontctx}{
      \oxextendctx{\oxvarctx_1}{
        \oxvarctxentry{\oxid}{\oxsitype}
      }
    }{\oxexpr}{\oxtunit}{
      \oxextendctx{\oxvarctx_1}{
        \oxvarctxentry{\oxid}{\oxsdtype}
      }
    }$ (from the premise of \oxname{T-ForArray}).

    To show $\oxtypjudge{\oxglobalctx}{\oxemptyctx}{\oxkontctx}{\oxvarctx_1}{
      \oxfor{\oxid}{\oxarr{\oxvalue_1 \oxdotsc \oxvalue_n}}{\oxexpr}
    }{\oxtunit}{\oxvarctx_1}$, we apply \Lemma{lemma:value-typing-output} to the
    derivation $\oxtypjudge{\oxglobalctx}{\oxemptyctx}{\oxkontctx}{\oxvarctx}
    {\oxarr{\oxvalue_0 \oxdotsc
        \oxvalue_n}}{\oxtarr{\oxsitype}{\oxnum}}{\oxvarctx_1}$ to get
    $\oxtypjudge{\oxglobalctx}{\oxemptyctx}{\oxkontctx}{\oxvarctx_1} {\oxarr{\oxvalue_0
        \oxdotsc \oxvalue_n}}{\oxtarr{\oxsitype}{\oxnum}}{\oxvarctx_1}$. Then,
    we rewrite the derivation (inverting and then reapply \oxname{T-Array}) to
    exclude $\oxvalue_0$ giving us
    $\oxtypjudge{\oxglobalctx}{\oxemptyctx}{\oxkontctx}{\oxvarctx_1} {\oxarr{\oxvalue_1
        \oxdotsc \oxvalue_n}}{\oxtarr{\oxsitype}{\oxnum - 1}}{\oxvarctx_1}$.
    Finally, we apply \oxname{T-ForArray} using this combined with
    $\oxtypjudge{\oxglobalctx}{\oxemptyctx}{\oxkontctx}{\oxvarctx_1}
    {\oxexpr}{\oxtunit}{\oxvarctx_1}$ (from the premise of \oxname{T-ForArray}).
  }

  %% new type is a subtype of the old type
  \oxpresline{
    \oxtunify[\oxcombine]{\oxemptyctx}{\oxkontctx}{\oxvarctx_1}
    {\oxtunit}{\oxtunit}{\oxvarctx_1}
  }{
    Immediate by \oxname{RR-Refl}.
  }

  %% new loan context is a sub-context of the old loan context
  \oxpresline{
    \exists \oxvarctx_o. \oxvarctx_1 \oxintersect \oxvarctx_o = \oxvarctx_1
  }{
    $\oxvarctx_o = \oxvarctx_1$
  }
\end{preservation}

\oxpreservationsubcaseheading{E-ForEmptyArray}{
  \oxvarctx
}

\begin{preservation}
  %% new store is valid
  \oxpresline{
    \oxstorevalidity{\oxglobalctx}{\oxvarctx}{\oxstore}
  }{
    Immediate from our premise.
  }

  %% kontctx values are still valid
  \oxpresline{
    \oxwfkontctx{\oxglobalctx}{\oxvarctx}{\overline{\oxvalue}}{\oxkontctx}
  }{
    Immediate from our premise.
  }

  %% new expression is well-typed
  \oxpresline{
    \oxtypjudge{\oxglobalctx}{\oxemptyctx}{\oxkontctx}{\oxvarctx}
    {\oxunit}{\oxtunit}{\oxvarctx}
  }{
    Immediate by \oxname{T-Unit}.
  }

  %% new type is a subtype of the old type
  \oxpresline{
    \oxtunify[\oxcombine]{\oxemptyctx}{\oxkontctx}{\oxvarctx}
    {\oxtunit}{\oxtunit}{\oxvarctx}
  }{
    Immediate by \oxname{RR-Refl}.
  }

  %% new loan context is a sub-context of the old loan context
  \oxpresline{
    \exists \oxvarctx_o. \oxvarctx \oxintersect \oxvarctx_o = \oxvarctx
  }{
    $\oxvarctx_o = \oxvarctx$
  }
\end{preservation}

\oxpreservationcasemulti{T-ForSlice}{\TForSlice}{
  \oxfor{\oxid}{\oxexpr_1}{\oxexpr_2}
}{
  \EForSlice \and \EForEmptySlice
}

\oxpreservationsubcaseheading{E-ForSlice}{
  \oxextendctx{\oxvarctx_1}{
    \oxvarctxentry{\oxid}{
      \oxtref{\oxcprov}{\oxmuta}{\oxsitype}
    }
  }
}

\begin{preservation}
  %% new store is valid
  \oxpresline{
    \oxstorevalidity{\oxglobalctx}{
      \oxextendctx{\oxvarctx_1}{
        \oxvarctxentry{\oxid}{
          \oxtref{\oxcprov}{\oxmuta}{\oxsitype}
        }
      }
    }{
      \oxextendctx{\oxstore}{
        \oxstoreentry{\oxid}{
          \oxindexptr{\oxreferent}{\oxnum_1}
        }
      }
    }
  }{
    Applying \Lemma{lemma:values-dont-change} to the typing derivation (from
    \oxname{T-ForSlice}) for $\oxexpr_1$ (which we know is a value from
    \oxname{E-ForSlice}) gives us
    $\oxsubtypectx{\oxglobalctx}{\oxemptyctx}{\oxvarctx}{\oxvarctx_1}$. Then,
    we can apply \Lemma{lemma:stack-validity-related-envs} to get
    $\oxstorevalidity{\oxglobalctx}{\oxvarctx_1}{\oxstore}$. Applying
    \Lemma{lemma:value-typing-output} to the derivation
    $\oxtypjudge{\oxglobalctx}{\oxemptyctx}{\oxkontctx}{\oxvarctx}
    {\oxsliceptr{\oxreferent}{i}{j}}{
      \oxtref{\oxcprov}{\oxmuta}{\oxtslice{\oxsitype}}
    }{\oxvarctx_1}$
    gives us $\oxtypjudge{\oxglobalctx}{\oxemptyctx}{\oxkontctx}{\oxvarctx_1}
    {\oxsliceptr{\oxreferent}{i}{j}}{
      \oxtref{\oxcprov}{\oxmuta}{\oxtslice{\oxsitype}}
    }{\oxvarctx_1}$.
    Then, using inversion (on \oxname{T-Pointer}), we get
    $\oxreferentvalidity{\oxglobalctx}{\oxvarctx_1}{
      \oxslice{\oxreferent}{i}{j}
    }{\oxxitype}$ (where $\oxxitype = \oxtarr{\oxsitype}{\oxnum}$ or
    $\oxtslice{\oxsitype}$) and
    $\oxloanpkg{\oxmuta}{\oxplace} \in \oxvarctx_1(\oxcprov)$ (where
    $\oxreferent = \oxreferentctx[\oxplace]$). We can invert
    \oxname{WF-RefSliceArray} or \oxname{WF-RefSliceSlice} (based on
    $\oxxitype$) for $\oxreferentvalidity{\oxglobalctx}{\oxvarctx_1}{
      \oxslice{\oxreferent}{i}{j}
    }{\oxxitype}$ and then apply \oxname{WF-RefIndexArray} or
    \oxname{WF-RefIndexSlice} appropriately to get
    $\oxreferentvalidity{\oxglobalctx}{\oxvarctx_1}{
      \oxindex{\oxreferent}{i}
    }{\oxxitype}$. We can then use
    \oxname{T-Pointer} to get
    $\oxtypjudge{\oxglobalctx}{\oxemptyctx}{\oxkontctx}{\oxvarctx_1}
    {\oxindexptr{\oxreferent}{\oxnum_1}}{
      \oxtref{\oxcprov}{\oxmuta}{\oxsitype}
    }{\oxvarctx_1}$. Finally, applying
    \Lemma{lemma:store-validity-extension} to
    $\oxstorevalidity{\oxglobalctx}{\oxvarctx_1}{\oxstore}$ and
    $\oxtypjudge{\oxglobalctx}{\oxemptyctx}{\oxkontctx}{\oxvarctx_1}
    {\oxindexptr{\oxreferent}{\oxnum_1}}{
      \oxtref{\oxcprov}{\oxmuta}{\oxsitype}
    }{\oxvarctx_1}$ gives us
    $\oxstorevalidity{\oxglobalctx}{
      \oxextendctx{\oxvarctx_1}{
        \oxvarctxentry{\oxid}{\oxtref{\oxcprov}{\oxmuta}{\oxsitype}}
      }
    }{
      \oxextendctx{\oxstore}{
        \oxstoreentry{\oxid}{\oxindexptr{\oxreferent}{\oxnum_1}}
      }
    }$.
  }

  %% kontctx values are still valid
  \oxpresline{
    \oxwfkontctx{\oxglobalctx}{
      \oxextendctx{\oxvarctx_1}{
        \oxvarctxentry{\oxid}{
          \oxtref{\oxcprov}{\oxmuta}{\oxsitype}
        }
      }
    }{\overline{\oxvalue}}{\oxkontctx}
  }{
    Inverting \oxname{WF-Temporaries} gives us $\forall i \in 1 \oxdots n. \;
    \oxtypjudge{\oxglobalctx}{\oxemptyctx}{\oxsitype_1 \oxdotsc \oxsitype_{i-1}}{\oxvarctx}{
      \oxvalue_i
    }{\oxsitype_i}{\oxvarctx}$.

    Applying \Lemma{lemma:values-dont-change} to the typing derivation (from
    \oxname{T-ForSlice}) for
    $\oxexpr_1$ (which we know is a value from \oxname{E-ForSlice}) gives us
    $\oxsubtypectx{\oxglobalctx}{\oxemptyctx}{\oxvarctx}{\oxvarctx_1}$.

    Thus, we can apply \Lemma{lemma:value-typing-related-envs} to each of the
    typing judgments to get $\forall i \in 1 \oxdots n. \;
    \oxtypjudge{\oxglobalctx}{\oxemptyctx}{\oxemptyctx}{\oxvarctx_1}{
      \oxvalue_i
    }{\oxsitype_i}{\oxvarctx_1}$.

    Applying \Lemma{lemma:values-well-typed-extension} to each of the value
    typing judgments and $\forall \oxcprov \in
    \oxfprovs{\oxtref{\oxprov}{\oxmuta}{\oxsitype}}. \;
    \oxnotreborrowed{\oxvarctx_1}{\oxcprov}$ (from the premise of
    \oxname{T-ForSlice}) gives us
    $\forall i \in 1 \oxdots n. \;
    \oxtypjudge{\oxglobalctx}{\oxemptyctx}{\oxemptyctx}{
      \oxextendctx{\oxvarctx_1}{
        \oxvarctxentry{\oxid}{
          \oxtref{\oxcprov}{\oxmuta}{\oxsitype}
        }
      }
    }{
      \oxvalue_i
    }{\oxsitype_i}{
      \oxextendctx{\oxvarctx_1}{
        \oxvarctxentry{\oxid}{
          \oxtref{\oxcprov}{\oxmuta}{\oxsitype}
        }
      }
    }$.
  }

  %% new expression is well-typed
  \oxpresline{
    \oxtypjudge{\oxglobalctx}{\oxemptyctx}{\oxkontctx}{
      \oxextendctx{\oxvarctx_1}{
        \oxvarctxentry{\oxid}{
          \oxtref{\oxcprov}{\oxmuta}{\oxsitype}
        }
      }
    }{\oxexpr^\prime}{\oxtunit}{\oxvarctx_1}
  }{
    We need to build a derivation for the expression $\oxseq{\oxshift{\oxexpr}}{
      \oxfor{\oxid}{
        \oxsliceptr{\oxreferent}{\oxnum_1^\prime}{\oxnum_2}
      }{\oxexpr}
    }$. The bottom of this derivation will be \oxname{T-Seq} which requires us
    to show that $\oxtypjudge{\oxglobalctx}{\oxemptyctx}{\oxkontctx}{
      \oxextendctx{\oxvarctx_1}{
        \oxvarctxentry{\oxid}{\oxsitype}
      }
    }{\oxshift{\oxexpr}}{\oxtunit}{\oxvarctx_1}$ and that
    $\oxtypjudge{\oxglobalctx}{\oxemptyctx}{\oxkontctx}{\oxvarctx_1}{
      \oxfor{\oxid}{
        \oxsliceptr{\oxreferent}{\oxnum_1^\prime}{\oxnum_2}
      }{\oxexpr}
    }{\oxtunit}{\oxvarctx_1}$.

    To show $\oxtypjudge{\oxglobalctx}{\oxemptyctx}{\oxkontctx}{
      \oxextendctx{\oxvarctx_1}{
        \oxvarctxentry{\oxid}{\oxsitype}
      }
    }{\oxshift{\oxexpr}}{\oxtunit}{\oxvarctx_1}$, we apply \oxname{T-Shift} to
    $\oxtypjudge{\oxglobalctx}{\oxemptyctx}{\oxkontctx}{
      \oxextendctx{\oxvarctx_1}{
        \oxvarctxentry{\oxid}{\oxsitype}
      }
    }{\oxexpr}{\oxtunit}{
      \oxextendctx{\oxvarctx_1}{
        \oxvarctxentry{\oxid}{\oxsdtype}
      }
    }$ (from the premise of \oxname{T-ForSlice}).

    To show $\oxtypjudge{\oxglobalctx}{\oxemptyctx}{\oxkontctx}{\oxvarctx_1}{
      \oxfor{\oxid}{\oxsliceptr{\oxreferent}{\oxnum_1^\prime}{\oxnum_2}}{\oxexpr}
    }{\oxtunit}{\oxvarctx_1}$, we apply
    \Lemma{lemma:value-typing-output} to the derivation
    $\oxtypjudge{\oxglobalctx}{\oxemptyctx}{\oxkontctx}{\oxvarctx}
    {\oxsliceptr{\oxreferent}{\oxnum_1}{\oxnum_2}}{
      \oxtref{\oxcprov}{\oxmuta}{\oxtslice{\oxsitype}}
    }{\oxvarctx_1}$
    to get $\oxtypjudge{\oxglobalctx}{\oxemptyctx}{\oxkontctx}{\oxvarctx_1}
    {\oxsliceptr{\oxreferent}{\oxnum_1}{\oxnum_2}}{
      \oxtref{\oxcprov}{\oxmuta}{\oxtslice{\oxsitype}}
    }{\oxvarctx_1}$. Then, we rewrite the derivation (inverting and then reapply
    \oxname{T-Pointer}) to increment $\oxnum_1$ to $\oxnum_1^\prime$ giving us
    $\oxtypjudge{\oxglobalctx}{\oxemptyctx}{\oxkontctx}{\oxvarctx_1}
    {\oxsliceptr{\oxreferent}{\oxnum_1^\prime}{\oxnum_2}}{
      \oxtref{\oxcprov}{\oxmuta}{\oxtslice{\oxsitype}}
    }{\oxvarctx_1}$.
    Finally, we apply \oxname{T-ForSlice} using this combined with
    $\oxtypjudge{\oxglobalctx}{\oxemptyctx}{\oxkontctx}{\oxvarctx_1}
    {\oxexpr}{\oxtunit}{\oxvarctx_1}$ (from the premise of \oxname{T-ForSlice}).
  }

  %% new type is a subtype of the old type
  \oxpresline{
    \oxtunify[\oxcombine]{\oxemptyctx}{\oxkontctx}{\oxvarctx_1}
    {\oxtunit}{\oxtunit}{\oxvarctx_1}
  }{
    Immediate by \oxname{RR-Refl}.
  }

  %% new loan context is a sub-context of the old loan context
  \oxpresline{
    \exists \oxvarctx_o. \oxvarctx_1 \oxintersect \oxvarctx_o = \oxvarctx_1
  }{
    $\oxvarctx_o = \oxvarctx_1$
  }
\end{preservation}

\oxpreservationsubcaseheading{E-ForEmptySlice}{
  \oxvarctx
}

\begin{preservation}
  %% new store is valid
  \oxpresline{
    \oxstorevalidity{\oxglobalctx}{\oxvarctx}{\oxstore}
  }{
    Immediate from our premise.
  }

  %% kontctx values are still valid
  \oxpresline{
    \oxwfkontctx{\oxglobalctx}{\oxvarctx}{\overline{\oxvalue}}{\oxkontctx}
  }{
    Immediate from our premise.
  }

  %% new expression is well-typed
  \oxpresline{
    \oxtypjudge{\oxglobalctx}{\oxemptyctx}{\oxkontctx}{\oxvarctx}
    {\oxunit}{\oxtunit}{\oxvarctx}
  }{
    Immediate by \oxname{T-Unit}.
  }

  %% new type is a subtype of the old type
  \oxpresline{
    \oxtunify[\oxcombine]{\oxemptyctx}{\oxkontctx}{\oxvarctx}
    {\oxtunit}{\oxtunit}{\oxvarctx}
  }{
    Immediate by \oxname{RR-Refl}.
  }

  %% new loan context is a sub-context of the old loan context
  \oxpresline{
    \exists \oxvarctx_o. \oxvarctx \oxintersect \oxvarctx_o = \oxvarctx
  }{
    $\oxvarctx_o = \oxvarctx$
  }
\end{preservation}

\oxpreservationcase[
% Let $\oxsitype_{\textsc{s}1} =
% \oxsubst{\oxsubst{\oxsitype_1}{\overline{\oxprov^\prime}}{\overline{\oxprov}}}
% {\overline{\oxsitype}}{\overline{\oxtvar}} \oxdotsc \oxsitype_{\textsc{s}n} =
% \oxsubst{\oxsubst{\oxsitype_n}{\overline{\oxprov^\prime}}{\overline{\oxprov}}}
% {\overline{\oxsitype}}{\overline{\oxtvar}}$.
]{T-AppFunction}{\TApp}{
  \oxapp{\oxexpr_f}{
    \overline{\oxenv} \oxcomma
    \overline{\oxprov^\prime} \oxcomma
    \overline{\oxsitype}
  }{
    \oxexpr_1 \oxdotsc \oxexpr_n
  }
}{\EAppFun}{
  \oxnewframe{\oxvarctx_b}{
    \oxvarctxentry{\oxid_1}{\delta(\oxsitype_1)} \oxdotsc
    \oxvarctxentry{\oxid_n}{\delta(\oxsitype_n)}
  }
}

\begin{preservation}
  %% new store is valid
  \oxpresline{
    \oxstorevalidity{\oxglobalctx}{
      \oxvarctx_i
    }{
      \oxstore^\prime
    }
  }{
    Applying \Lemma{lemma:values-dont-change} to the derivation for $\oxvalue_f$
    gives us $\oxsubtypectx{\oxglobalctx}{\oxemptyctx}{\oxvarctx}{\oxvarctx_0}$.
    Then, applying \Lemma{lemma:values-dont-change} to the derivations for
    every $\oxvalue_i$ gives us $\forall i \in \oxset{1 \oxdots n}. \;
    \oxsubtypectx{\oxglobalctx}{\oxemptyctx}{\oxvarctx_{i-1}}{\oxvarctx_i}$.

    Inversion on $\oxrunifymany{\oxemptyctx}{\oxkontctx}{\oxvarctx_n}{
      \oxsubstmany{\oxabsprov_2}{\oxprov}{\oxabsprov}
    }{
      \oxsubstmany{\oxabsprov_1}{\oxprov}{\oxabsprov}
    }{\oxvarctx_b}$ gives us a sequence of outlives relations with
    intermediate contexts. Applying  \Lemma{lemma:outlives-related} to each of
    them and then combining the result by transitivity gives us
    $\oxsubtypectx{\oxglobalctx}{\oxemptyctx}{\oxvarctx_n}{\oxvarctx_b}$.
    Combining both by transitivity, we have
    $\oxsubtypectx{\oxglobalctx}{\oxemptyctx}{\oxvarctx}{\oxvarctx_b}$.

    We can then apply \Lemma{lemma:stack-validity-related-envs} with
    $\oxstorevalidity{\oxglobalctx}{\oxvarctx}{\oxstore}$ to get
    $\oxstorevalidity{\oxglobalctx}{\oxvarctx_b}{\oxstore}$.

    We can apply \oxname{WF-StackFrame} to get
    $\oxstorevalidity{\oxglobalctx}{
      \oxnewframe{\oxvarctx_b}{\oxemptyctx}
    }{
      \oxnewframe{\oxstore}{\oxemptyctx}
    }$. Then, we repeatedly apply \Lemma{lemma:store-validity-extension} to
    the derivations for the arguments ($\oxvalue_1 \oxdots \oxvalue_n$) to get
    $\oxstorevalidity{\oxglobalctx}{
      \oxnewframe{\oxvarctx}{
        \oxvarctxentry{\oxid_1}{\delta(\oxsitype_1)} \oxdotsc
        \oxvarctxentry{\oxid_n}{\delta(\oxsitype_n)}
      }
    }{
      \oxnewframe{\oxstore}{
        \oxstoreentry{\oxid_1}{\oxvalue_1} \oxdotsc
        \oxstoreentry{\oxid_n}{\oxvalue_n}
      }
    }$.
  }

  %% kontctx values are still valid
  \oxpresline{
    \oxwfkontctx{\oxglobalctx}{\oxvarctx_i}{\overline{\oxvalue}}{\oxkontctx}
  }{
    Inverting \oxname{WF-Temporaries} gives us $\forall i \in 1 \oxdots n. \;
    \oxtypjudge{\oxglobalctx}{\oxemptyctx}{\oxsitype_1 \oxdotsc \oxsitype_{i-1}}{\oxvarctx}{
      \oxvalue_i
    }{\oxsitype_i}{\oxvarctx}$.

    Applying \Lemma{lemma:values-dont-change} to the derivation for $\oxvalue_f$
    gives us $\oxsubtypectx{\oxglobalctx}{\oxemptyctx}{\oxvarctx}{\oxvarctx_0}$.
    Then, applying \Lemma{lemma:values-dont-change} to the derivations for
    every $\oxvalue_i$ gives us $\forall i \in \oxset{1 \oxdots n}. \;
    \oxsubtypectx{\oxglobalctx}{\oxemptyctx}{\oxvarctx_{i-1}}{\oxvarctx_i}$.

    Inversion on $\oxrunifymany{\oxemptyctx}{\oxkontctx}{\oxvarctx_n}{
      \oxsubstmany{\oxabsprov_2}{\oxprov}{\oxabsprov}
    }{
      \oxsubstmany{\oxabsprov_1}{\oxprov}{\oxabsprov}
    }{\oxvarctx_b}$ gives us a sequence of outlives relations with
    intermediate contexts. Applying  \Lemma{lemma:outlives-related} to each of
    them and then combining the result by transitivity gives us
    $\oxsubtypectx{\oxglobalctx}{\oxemptyctx}{\oxvarctx_n}{\oxvarctx_b}$.
    Combining both by transitivity, we have
    $\oxsubtypectx{\oxglobalctx}{\oxemptyctx}{\oxvarctx}{\oxvarctx_b}$.

    Thus, we can apply \Lemma{lemma:value-typing-related-envs} to each of the
    typing judgments to get $\forall i \in 1 \oxdots n. \;
    \oxtypjudge{\oxglobalctx}{\oxemptyctx}{\oxemptyctx}{\oxvarctx_b}{
      \oxvalue_i
    }{\oxsitype_i}{\oxvarctx_b}$.

    We also immediately have that adding a new empty frame leaves all the typing
    derivations good since it doesn't actually bind any new names on its own.
    Thus, we have $\forall i \in 1 \oxdots n. \;
    \oxtypjudge{\oxglobalctx}{\oxemptyctx}{\oxemptyctx}{
      \oxnewframe{\oxvarctx_b}{\oxemptyctx}
    }{
      \oxvalue_i
    }{\oxsitype_i}{
      \oxnewframe{\oxvarctx_b}{\oxemptyctx}
    }$.

    Repeatedly applying \Lemma{lemma:values-well-typed-extension} to each of the value
    typing judgments gives us
    $\forall i \in 1 \oxdots n. \;
    \oxtypjudge{\oxglobalctx}{\oxemptyctx}{\oxemptyctx}{
      \oxnewframe{\oxvarctx_b}{
        \oxvarctxentry{\oxid_1}{\delta(\oxsitype_1)} \oxdotsc
        \oxvarctxentry{\oxid_n}{\delta(\oxsitype_n)}
      }
    }{
      \oxvalue_i
    }{\oxsitype_i}{
      \oxnewframe{\oxvarctx_b}{
        \oxvarctxentry{\oxid_1}{\delta(\oxsitype_1)} \oxdotsc
        \oxvarctxentry{\oxid_n}{\delta(\oxsitype_n)}
      }
    }$.
  }

  %% new expression is well-typed
  \oxpresline{
    \oxtypjudge{\oxglobalctx}{\oxemptyctx}{\oxkontctx}{\oxvarctx_i}{
      \oxexpr^\prime
    }{\delta(\oxsitype_f)}{\oxvarctx_b}
  }{
    Applying \Lemma{lemma:function-defns-self-contained} with
    $\oxctxswellformed{\oxglobalctx}{\oxemptyctx}{\oxvarctx_b}{\oxkontctx}$ and
    $\oxlookup{\oxglobalctx}{\oxfnname}{
      \oxfuncdef{\oxfnname}{
        \overline{\oxenvvar} \oxcomma \overline{\oxabsprov} \oxcomma \overline{\oxtvar}
      }{
        \oxascribe{\oxid_1}{\oxsitype_1} \oxdotsc
        \oxascribe{\oxid_n}{\oxsitype_n}
      }{\oxsitype_r}{\overline{\oxabsprov_1: \oxabsprov_2}}{\oxexpr}
    }$ (from the premise of on \oxname{E-AppFunction}) gives us
    $\oxtypjudge{\oxglobalctx}{
      \overline{\oxtvarctxentry{\oxenvvar}{\oxkenv}} \oxcomma
      \overline{\oxtvarctxentry{\oxabsprov}{\oxkprov}} \oxcomma
      \overline{\oxabsprov_1 \oxoutlives \oxabsprov_2} \oxcomma
      \overline{\oxtvarctxentry{\oxtvar}{\oxktype}}
    }{\oxkontctx}{
      \oxnewframe{\oxvarctx_b}{
        \oxvarctxentry{\oxid_1}{\oxsitype_1} \oxdotsc
        \oxvarctxentry{\oxid_n}{\oxsitype_n}
      }
    }{
      \oxframed{\oxexpr}
    }{
      \oxsitype_f
    }{\oxvarctx_b}$.

    We then repeatedly apply \Lemma{lemma:substitution} for all of our type,
    region, and environment variables to get
    $\oxtypjudge{\oxglobalctx}{\oxemptyctx}{\oxkontctx}{
      \oxnewframe{\oxvarctx_b}{
        \oxvarctxentry{\oxid_1}{\delta(\oxsitype_1)} \oxdotsc
        \oxvarctxentry{\oxid_n}{\delta(\oxsitype_n)}
      }
    }{
      \oxframed{\oxexpr}
    }{\delta(\oxsitype_f)}{\oxvarctx_b}$.
  }

  %% new type is a subtype of the old type
  \oxpresline{
    \oxtunify[\oxcombine]{\oxemptyctx}{\oxkontctx}{\oxvarctx_b}{
      \oxtype^\prime
    }{
      \oxtype^\prime
    }{\oxvarctx_b}
  }{
    Immediate by \oxname{RR-Refl}.
  }

  %% new loan context is a sub-context of the old loan context
  \oxpresline{
    \exists \oxvarctx_o. \oxvarctx_b \oxintersect \oxvarctx_o = \oxvarctx_b
  }{
    $\oxvarctx_o = \oxvarctx_b$
  }
\end{preservation}

\oxpreservationcase{T-AppClosure}{\TAppClosure}{
  \oxapply{\oxexpr_f}{
    \oxexpr_1 \oxdotsc \oxexpr_n
  }
}{\EApp}{
  \oxnewframe{\oxvarctx^\prime_n}{
    \oxextendctx{\oxframe_c}{
      \oxvarctxentry{\oxid_1}{\oxsitype_1} \oxdotsc
      \oxvarctxentry{\oxid_n}{\oxsitype_n}
    }
  }
}

\begin{preservation}
  %% new store is valid
  \oxpresline{
    \oxstorevalidity{\oxglobalctx}{
      \oxvarctx_i
    }{
      \oxstore^\prime
    }
  }{
    Applying \Lemma{lemma:values-dont-change} to the derivation for $\oxvalue_f$
    gives us $\oxsubtypectx{\oxglobalctx}{\oxemptyctx}{\oxvarctx}{\oxvarctx_0}$.
    We can apply \Lemma{lemma:stack-validity-related-envs} with
    $\oxstorevalidity{\oxglobalctx}{\oxvarctx}{\oxstore}$ to get
    $\oxstorevalidity{\oxglobalctx}{\oxvarctx_0}{\oxstore}$.

    Then, for each pair of value typing judgment and region rewriting, we follow
    the same pattern we will describe in terms of an arbitrary index $i$. First,
    we apply \Lemma{lemma:values-dont-change} to the derivation for $\oxvalue_i$
    to get $\oxsubtypectx{\oxglobalctx}{\oxemptyctx}
    {\oxvarctx_{i-1}}{\oxvarctx_i}$. Then, applying
    \Lemma{lemma:stack-validity-related-envs} gives us
    $\oxstorevalidity{\oxglobalctx}{\oxvarctx_i}{\oxstore}$. Then, we can apply
    \Lemma{lemma:stack-validity-region-rewriting} to
    $\oxtunify[\oxcombineunrest]{\oxemptyctx}{\oxkontctx}{\oxvarctx_i}
    {\oxsitype_{i^\prime}}{\oxsitype_i}{\oxvarctx_i^\prime}$ to get
    $\oxstorevalidity{\oxglobalctx}{\oxvarctx_i^\prime}{\oxstore}$. Going through
    this for all the values in sequence gives us
    $\oxstorevalidity{\oxglobalctx}{\oxvarctx_n^\prime}{\oxstore}$

    Then, inversion of \oxname{T-ClosureValue} for the typing derivation for
    $\oxvalue_f$ gives us
    $\oxsubstorevalidity{\oxglobalctx}{\oxvarctx}{\oxstackframe_c}{\oxframe_c}$.
    We can then invert \oxname{WF-Frame} here to get
    $\oxdomain{\oxstackframe} = \oxdomain{\oxframe_c}|_\oxid$
    $\forall \oxid \in \oxdomain{\oxstackframe}. \;
    \oxtypjudge{\oxglobalctx}{\oxemptyctx}{\oxkontctx}{\oxnewframe{\oxvarctx}{\oxframe_c}}
    {\oxstackframe(\oxid)}{\oxframe_c(\oxid)}{\oxnewframe{\oxvarctx}{\oxframe_c}}$
    which we can then use with $\oxstorevalidity{\oxglobalctx}{\oxvarctx_n}{\oxstore}$
    in \oxname{WF-StackFrame} to get
    $\oxstorevalidity{\oxglobalctx}{
      \oxnewframe{\oxvarctx_n^{\prime}}{\oxframe_c}
    }{
      \oxnewframe{\oxstore}{\oxstackframe_c}
    }$.

    Finally, we repeatedly apply \Lemma{lemma:store-validity-extension} to
    the derivations for the arguments ($\oxvalue_1 \oxdots \oxvalue_n$) to get
    $\oxstorevalidity{\oxglobalctx}{
      \oxvarctx_i
    }{
      \oxstore^\prime
    }$.
  }

  %% kontctx values are still valid
  \oxpresline{
    \oxwfkontctx{\oxglobalctx}{\oxvarctx_i}{\overline{\oxvalue^\prime}}{\oxkontctx}
  }{
    Inverting \oxname{WF-Temporaries} gives us $\forall j \in 1 \oxdots n. \;
    \oxtypjudge{\oxglobalctx}{\oxemptyctx}{\oxsitype_1 \oxdotsc \oxsitype_{i-1}}{\oxvarctx}{
      \oxvalue^\prime_j
    }{\oxsitype_j}{\oxvarctx}$.

    Applying \Lemma{lemma:values-dont-change} to the derivation for $\oxvalue_f$
    gives us $\oxsubtypectx{\oxglobalctx}{\oxemptyctx}{\oxvarctx}{\oxvarctx_0}$.
    For each of $\forall j \in 1 \oxdots n. \;
    \oxtypjudge{\oxglobalctx}{\oxemptyctx}{\oxemptyctx}{\oxvarctx}{
      \oxvalue^\prime_j
    }{\oxsitype_i}{\oxvarctx}$, we apply \Lemma{lemma:value-typing-related-envs}
    to get $\forall j \in 1 \oxdots n. \;
    \oxtypjudge{\oxglobalctx}{\oxemptyctx}{\oxemptyctx}{\oxvarctx_0}{
      \oxvalue^\prime_j
    }{\oxsitype_j}{\oxvarctx_0}$

    From here, we have an alternating pattern of lemmas to apply to deal with
    the interleaved value typing and then region rewriting in the premise of
    \oxname{T-AppClosure}. For each value $\oxvalue_i$ in \oxname{T-AppClosure},
    we'll first apply \Lemma{lemma:values-dont-change} to get
    $\oxsubtypectx{\oxglobalctx}{\oxemptyctx}{\oxvarctx_{i-1}}{\oxvarctx_i}$.
    Then, we'll apply \Lemma{lemma:value-typing-related-envs} to
    $\forall j \in 1 \oxdots n. \;
    \oxtypjudge{\oxglobalctx}{\oxemptyctx}{\oxemptyctx}{\oxvarctx_{i-1}}{
      \oxvalue^\prime_j
    }{\oxsitype_j}{\oxvarctx_{i-1}}$ to get $\forall j \in 1 \oxdots n. \;
    \oxtypjudge{\oxglobalctx}{\oxemptyctx}{\oxemptyctx}{\oxvarctx_i}{
      \oxvalue^\prime_j
    }{\oxsitype_j}{\oxvarctx_i}$. Then, we apply
    \Lemma{lemma:subtyping-value-typing} to each of these along with the
    rewriting $\oxtunify[\oxcombineunrest]{\oxemptyctx}{\oxkontctx}{\oxvarctx_i}
    {\oxsitype_{i^\prime}}{\oxsitype_i}{\oxvarctx_i^\prime}$ to get
    $\forall j \in 1 \oxdots n. \;
    \oxtypjudge{\oxglobalctx}{\oxemptyctx}{\oxemptyctx}{\oxvarctx_i^\prime}{
      \oxvalue^\prime_j
    }{\oxsitype_j}{\oxvarctx_i^\prime}$. After going through this for each value
    $\oxvalue_i$, we have $\forall j \in 1 \oxdots n. \;
    \oxtypjudge{\oxglobalctx}{\oxemptyctx}{\oxemptyctx}{\oxvarctx_n^\prime}{
      \oxvalue^\prime_j
    }{\oxsitype_j}{\oxvarctx_n^\prime}$.

    We also immediately have that adding a new empty frame leaves all the typing
    derivations good since it doesn't actually bind any new names on its own.
    Thus, we have $\forall i \in 1 \oxdots n. \;
    \oxtypjudge{\oxglobalctx}{\oxemptyctx}{\oxemptyctx}{
      \oxnewframe{\oxvarctx_n^\prime}{\oxemptyctx}
    }{
      \oxvalue_i
    }{\oxsitype_i}{
      \oxnewframe{\oxvarctx_n^\prime}{\oxemptyctx}
    }$.

    Repeatedly applying \Lemma{lemma:values-well-typed-extension} to each of the
    value typing judgments and each of $\forall i \in \oxset{ 1 \oxdots n }. \;
    \forall \oxcprov \in \oxfprovs{\oxsitype_i}. \;
    \oxnotreborrowed{\oxvarctx_1}{\oxcprov}$ (from the premise of
    \oxname{T-App}) gives us
    $\forall i \in 1 \oxdots n. \;
    \oxtypjudge{\oxglobalctx}{\oxemptyctx}{\oxemptyctx}{
      \oxnewframe{\oxvarctx_n^\prime}{
        \oxextendctx{\oxframe_c}{
          \oxvarctxentry{\oxid_1}{\oxsitype_1} \oxdotsc
          \oxvarctxentry{\oxid_n}{\oxsitype_n}
        }
      }
    }{
      \oxvalue_i
    }{\oxsitype_i}{
      \oxnewframe{\oxvarctx_n^\prime}{
        \oxextendctx{\oxframe_c}{
          \oxvarctxentry{\oxid_1}{\oxsitype_1} \oxdotsc
          \oxvarctxentry{\oxid_n}{\oxsitype_n}
        }
      }
    }$.
  }
\end{preservation} %% breaking for pagebreak
\begin{preservation}
  %% new expression is well-typed
  \oxpresline{
    \oxtypjudge{\oxglobalctx}{\oxemptyctx}{\oxkontctx}{\oxvarctx_i}{
      \oxexpr^\prime
    }{\oxsitype_f}{\oxvarctx^\prime_n}
  }{
    Consider first the typing derivation for $\oxvalue_f$ of
    $\oxtypjudge{\oxglobalctx}{\oxemptyctx}{\oxkontctx}{\oxvarctx}{
      \oxvalue_f
    }{
      \oxtfun{}{
        \oxsitype_1 \oxdotsc \oxsitype_n
      }{\oxsitype_f}{\oxenv_c}
    }{\oxvarctx_0^\prime}$.

    Applying \Lemma{lemma:value-typing-output} gives us
    $\oxtypjudge{\oxglobalctx}{\oxemptyctx}{\oxkontctx}{\oxvarctx_0^\prime}{
      \oxvalue_f
    }{
      \oxtfun{}{
        \oxsitype_1 \oxdotsc \oxsitype_n
      }{\oxsitype_f}{\oxenv_c}
    }{\oxvarctx_0^\prime}$.

    Then, as in the previous cases, we can apply
    \Lemma{lemma:values-dont-change}, \Lemma{lemma:value-typing-related-envs},
    and \Lemma{lemma:subtyping-value-typing} in sequence starting with the above
    typing derivation with each of the typing derivations for the arguments
    $\oxvalue_i$ and their subsequent corresponding region rewriting derivation.
    At the end, this gives us
    $\oxtypjudge{\oxglobalctx}{\oxemptyctx}{\oxkontctx}{\oxvarctx_n^\prime}{
      \oxvalue_f
    }{
      \oxtfun{}{
        \oxsitype_1 \oxdotsc \oxsitype_n
      }{\oxsitype_f}{\oxenv_c}
    }{\oxvarctx_n^\prime}$

    Then, inversion on \oxname{T-ClosureValue} on this typing derivation for
    $\oxvalue_f$ gives us $\oxfvars{\oxexpr} \setminus \overline{\oxid} =
    \overline{\oxid_f} = \oxdomain{\oxframe_c}|_\oxid$,
    $\overline{\oxcprov} = \overline{\oxfprovs{\oxvarctx(\oxid_f)}} \oxcomma
    \oxfprovs{\oxexpr} = \oxdomain{\oxframe_c}|_\oxcprov$, and
    $\oxtypjudge{\oxglobalctx}{\oxemptyctx}{\oxkontctx}{
      \oxnewframe{\oxvarctx_n^\prime}{
        \oxframe_c \oxcomma
        \oxvarctxentry{\oxid_1}{\oxsitype_1} \oxdotsc
        \oxvarctxentry{\oxid_n}{\oxsitype_n}
      }
    }{\oxexpr}{\oxsitype_r}{\oxnewframe{\oxvarctx_n^\prime}{\oxframe}}$. We can
    then apply \oxname{T-Framed} to get
    $\oxtypjudge{\oxglobalctx}{\oxemptyctx}{\oxkontctx}{
      \oxnewframe{\oxvarctx_n}{
        \oxextendctx{\oxframe_c}{
          \oxvarctxentry{\oxid_1}{\oxsitype_1} \oxdotsc
          \oxvarctxentry{\oxid_n}{\oxsitype_n}
        }
      }
    }{
      \oxframed{\oxexpr}
    }{\oxsitype_f}{\oxvarctx_n}$.
  }

  %% new type is a subtype of the old type
  \oxpresline{
    \oxtunify[\oxcombine]{\oxemptyctx}{\oxkontctx}{\oxvarctx^\prime_n}{
      \oxsitype_f
    }{
      \oxsitype_f
    }{\oxvarctx^\prime_n}
  }{
    Immediate by \oxname{RR-Refl}.
  }

  %% new loan context is a sub-context of the old loan context
  \oxpresline{
    \exists \oxvarctx_o. \oxvarctx^\prime_n \oxintersect \oxvarctx_o =
    \oxvarctx^\prime_n
  }{
    $\oxvarctx_o = \oxvarctx^\prime_n$
  }
\end{preservation}

\oxpreservationcase{T-Function}{\TFunction}{
  \oxtfun{\overline{\oxprov} \oxcomma \overline{\oxtvar}}{
    \oxsitype_1 \oxdotsc \oxsitype_n
  }{\oxsitype_r}{}
}{\EFunction}{
  \oxvarctx
}

\begin{preservation}
  %% new store is valid
  \oxpresline{\oxstorevalidity{\oxglobalctx}{\oxvarctx}{\oxstore}}{
    Immediate from our premise.
  }

  %% kontctx values are still valid
  \oxpresline{
    \oxwfkontctx{\oxglobalctx}{\oxvarctx}{\overline{\oxvalue}}{\oxkontctx}
  }{
    Immediate from our premise.
  }

  %% new expression is well-typed
  \oxpresline{
    \oxtypjudge{\oxglobalctx}{\oxemptyctx}{\oxkontctx}{\oxvarctx}
    {\oxexpr^\prime}{\oxtype^\prime}{\oxvarctx}
  }{
    By \oxname{T-ClosureValue} since $\oxstore_c$ is empty, and we know that the
    body itself is well-typed as a consequence of inversion on
    \oxname{WF-FunctionDefinition} for \oxfnname.
  }

  %% new type is a subtype of the old type
  \oxpresline{
    \oxtunify[\oxcombine]{\oxemptyctx}{\oxkontctx}{\oxvarctx}
    {\oxtype^\prime}{\oxtype}{\oxvarctx}
  }{
    Immediate by \oxname{RR-Refl}.
  }

  %% new loan context is a sub-context of the old loan context
  \oxpresline{
    \exists \oxvarctx_o. \oxvarctx \oxintersect \oxvarctx_o = \oxvarctx
  }{
    $\oxvarctx_o = \oxvarctx$
  }
\end{preservation}


\oxpreservationcase{T-Closure}{\TClosure}{
  \oxfunction{
    \overline{\oxenvvar} \oxcomma \overline{\oxprov} \oxcomma \overline{\oxtvar}
  }{
    \oxascribe{\oxid_1}{\oxsitype_1} \oxdotsc
    \oxascribe{\oxid_n}{\oxsitype_n}
  }{\oxsitype_r}{
    \oxexpr
  }
}{\EClosure}{
  \oxtupdatemany{\oxvarctx}{\oxid_{nc}}{\oxtuninit{\oxvarctx(\oxid_{nc})}}
}

\begin{preservation}
  %% new store is valid
  \oxpresline{
    \oxstorevalidity{\oxglobalctx}{
      \oxtupdatemany{\oxvarctx}{\oxid_{nc}}{\oxtuninit{\oxvarctx(\oxid_{nc})}}
    }{
      \oxupdate{\oxstore}{
        \overline{\oxstoreentry{\oxid_{nc}}{\oxuninit}}
      }
    }
  }{
    Compared to $\oxvarctx$, we know that
    $\oxtupdatemany{\oxvarctx}{\oxid_{nc}}{\oxtuninit{\oxvarctx(\oxid_{nc})}}$
    has more things marked dead and no other changes. Thus, we can apply
    \oxname{R-Env} to get $\oxsubtypectx{\oxglobalctx}{\oxemptyctx}{\oxvarctx}{
      \oxtupdatemany{\oxvarctx}{\oxid_{nc}}{\oxtuninit{\oxvarctx(\oxid_{nc})}}
    }$. Then, we can apply \Lemma{lemma:stack-validity-related-envs} to get
    $\oxstorevalidity{\oxglobalctx}{
      \oxtupdatemany{\oxvarctx}{\oxid_{nc}}{\oxtuninit{\oxvarctx(\oxid_{nc})}}
    }{\oxstore}$. Since $\oxuninit$ is good at every dead type $\oxsdtype$ by
    \oxname{T-Dead}, we can then build a new derivation using that rule instead
    for every $\oxid_{nc}$ that is now at a dead type. This gives us
    $\oxstorevalidity{\oxglobalctx}{
      \oxtupdatemany{\oxvarctx}{\oxid_{nc}}{\oxtuninit{\oxvarctx(\oxid_{nc})}}
    }{
      \oxupdate{\oxstore}{
        \overline{\oxstoreentry{\oxid_{nc}}{\oxuninit}}
      }
    }$.
  }

  %% kontctx values are still valid
  \oxpresline{
    \oxwfkontctx{\oxglobalctx}{
      \oxtupdatemany{\oxvarctx}{\oxid_{nc}}{\oxtuninit{\oxvarctx(\oxid_{nc})}}
    }{\overline{\oxvalue}}{\oxkontctx}
  }{
    Inverting \oxname{WF-Temporaries} gives us $\forall i \in 1 \oxdots n. \;
    \oxtypjudge{\oxglobalctx}{\oxemptyctx}{\oxsitype_1 \oxdotsc \oxsitype_{i-1}}{\oxvarctx}{
      \oxvalue_i
    }{\oxsitype_i}{\oxvarctx}$. By \oxname{R-Env}, we have that
    \oxsubtypectx{\oxglobalctx}{\oxemptyctx}{\oxvarctx}{
      \oxtupdatemany{\oxvarctx}{\oxid_{nc}}{\oxtuninit{\oxvarctx(\oxid_{nc})}}
    }. Thus, we can apply \Lemma{lemma:value-typing-related-envs} to each of the
    typing judgments and then apply \oxname{WF-Temporaries} to get
    $\oxwfkontctx{\oxglobalctx}{
      \oxtupdatemany{\oxvarctx}{\oxid_{nc}}{\oxtuninit{\oxvarctx(\oxid_{nc})}}
    }{\overline{\oxvalue}}{\oxkontctx}$.
  }

  %% new expression is well-typed
  \oxpresline{
    \oxtypjudge{\oxglobalctx}{\oxemptyctx}{\oxkontctx}{\oxvarctx_i}
    {\oxexpr^\prime}{\oxtype^\prime}{\oxvarctx_i}
  }{
    Immediate by inversion of \oxname{T-Closure} and application of
    \oxname{T-ClosureValue}. Note that they have identical premises.
  }

  %% new type is a subtype of the old type
  \oxpresline{
    \oxtunify[\oxcombine]{\oxemptyctx}{\oxkontctx}{\oxvarctx_i}
    {\oxtype^\prime}{\oxtype^\prime}{\oxvarctx_i}
  }{
    Immediate by \oxname{RR-Refl}.
  }

  %% new loan context is a sub-context of the old loan context
  \oxpresline{
    \exists \oxvarctx_o. \oxvarctx_i \oxintersect \oxvarctx_o = \oxvarctx_i
  }{
    $\oxvarctx_o = \oxvarctx_i$
  }
\end{preservation}

\oxpreservationcase{T-Shift}{\TShift}{
  \oxshift{\oxexpr_i}
}{\EShift}{\oxvarctx^\prime}

\begin{preservation}
  %% new store is valid
  \oxpresline{
    \oxstorevalidity{\oxglobalctx}{\oxvarctx^\prime}{\oxstore}
  }{
    By inversion of \oxname{WF-StackFrame} on $\oxstorevalidity{\oxglobalctx}{
      \oxextendctx{\oxvarctx^\prime}{\oxvarctxentry{\oxid}{\oxsdtype}}
    }{\oxextendctx{\oxstore}{\oxstoreentry{\oxid}{\oxvalue^\prime}}}$, we get
    $\oxstorevalidity{\oxglobalctx}{\oxvarctx_i}{\oxstore}$,
    $\oxdomain{
      \oxextendctx{\oxstackframe}{
        \oxstoreentry{\oxid}{\oxvalue^\prime}
      }
    } = \oxdomain{
      \oxextendctx{\oxframe}{
        \oxvarctxentry{\oxid}{\oxsdtype}
      }
    }|_\oxid$, and
    $\forall \oxid \in \oxdomain{
      \oxnewframe{\oxstore}{
        \oxextendctx{\oxstackframe}{
          \oxstoreentry{\oxid}{\oxvalue^\prime}
        }
      }
    }. \;
    \oxtypjudge{\oxglobalctx}{\oxemptyctx}{\oxkontctx}{
      \oxnewframe{\oxvarctx_i}{
        \oxextendctx{\oxframe}{
          \oxvarctxentry{\oxid}{\oxsdtype}
        }
      }
    }{
      (\oxnewframe{\oxstore}{
        \oxextendctx{\oxstackframe}{
          \oxstoreentry{\oxid}{\oxvalue^\prime}
        }
      })(\oxid)
    }{(\oxnewframe{\oxvarctx_i}{
        \oxextendctx{\oxframe}{
          \oxvarctxentry{\oxid}{\oxsdtype}
        }
      })(\oxid)}{
      \oxnewframe{\oxvarctx_i}{
        \oxextendctx{\oxframe}{
          \oxvarctxentry{\oxid}{\oxsdtype}
        }
      }
    }$. Note that $\oxnewframe{\oxvarctx_i}{\oxframe} = \oxvarctx^\prime$.
    We can then immediately see that the above implies
    $\oxdomain{\oxstackframe} = \oxdomain{\oxframe}|_\oxid$ and
    $\forall \oxid \in \oxdomain{
      \oxnewframe{\oxstore}{
        \oxextendctx{\oxstackframe}{
          \oxstoreentry{\oxid}{\oxvalue^\prime}
        }
      }
    }. \;
    \oxtypjudge{\oxglobalctx}{\oxemptyctx}{\oxkontctx}{
      \oxnewframe{\oxvarctx_i}{\oxframe}
    }{
      (\oxnewframe{\oxstore}{\oxstackframe})(\oxid)
    }{(\oxnewframe{\oxvarctx_i}{\oxframe})(\oxid)}{
      \oxnewframe{\oxvarctx_i}{\oxframe}
    }$. Thus, we can apply \oxname{WF-StackFrame} to get
    $\oxstorevalidity{\oxglobalctx}{
      \oxnewframe{\oxvarctx_i}{\oxstackframe}
    }{\oxstore}$ which can be rewritten as
    $\oxstorevalidity{\oxglobalctx}{\oxvarctx^\prime}{\oxstore}$.
  }

  %% kontctx values are still valid
  \oxpresline{
    \oxwfkontctx{\oxglobalctx}{\oxvarctx^\prime}{\overline{\oxvalue}}{\oxkontctx}
  }{
    Inverting \oxname{WF-Temporaries} gives us $\forall i \in 1 \oxdots n. \;
    \oxtypjudge{\oxglobalctx}{\oxemptyctx}{\oxsitype_1 \oxdotsc \oxsitype_{i-1}}{\oxvarctx}{
      \oxvalue_i
    }{\oxsitype_i}{\oxvarctx}$.

    Applying \Lemma{lemma:values-dont-change} to the typing derivation (from
    \oxname{T-Shift}) for
    $\oxexpr$ (which we know is a value from \oxname{E-Shift}) gives us
    $\oxsubtypectx{\oxglobalctx}{\oxemptyctx}{\oxvarctx}{\oxvarctx^\prime}$.

    Thus, we can apply \Lemma{lemma:value-typing-related-envs} to each of the
    typing judgments to get $\forall i \in 1 \oxdots n. \;
    \oxtypjudge{\oxglobalctx}{\oxemptyctx}{\oxsitype_1 \oxdotsc \oxsitype_{i-1}}{\oxvarctx^\prime}{
      \oxvalue_i
    }{\oxsitype_i}{\oxvarctx^\prime}$.

    Finally, applying \oxname{WF-Temporaries} gives us
    $\oxwfkontctx{\oxglobalctx}{\oxvarctx^\prime}{\overline{\oxvalue}}{\oxkontctx}$.
  }

  %% new expression is well-typed
  \oxpresline{
    \oxtypjudge{\oxglobalctx}{\oxemptyctx}{\oxkontctx}{\oxvarctx^\prime}{\oxvalue}
    {\oxsitype}{\oxvarctx^\prime}
  }{
    We know from \oxname{E-Shift} that $\oxexpr$ is a value $\oxvalue$. Thus,
    we can apply \Lemma{lemma:value-typing-output} to
    $\oxtypjudge{\oxglobalctx}{\oxemptyctx}{\oxkontctx}{\oxvarctx}{\oxvalue}{\oxsitype}{
      \oxextendctx{\oxvarctx^\prime}{
        \oxvarctxentry{\oxid}{\oxsdtype}
      }
    }$ to get $\oxtypjudge{\oxglobalctx}{\oxemptyctx}{\oxkontctx}{
      \oxextendctx{\oxvarctx^\prime}{
        \oxvarctxentry{\oxid}{\oxsdtype}
      }
    }{\oxvalue}{\oxsitype}{
      \oxextendctx{\oxvarctx^\prime}{
        \oxvarctxentry{\oxid}{\oxsdtype}
      }
    }$.

    We now wish to show that
    $\oxtypjudge{\oxglobalctx}{\oxemptyctx}{\oxkontctx}{\oxvarctx^\prime}{\oxvalue}
    {\oxsitype}{\oxvarctx^\prime}$. By inspecting the grammar of values and
    their typing rules, we know that the only values who depend on the
    context are pointers and closure values. But by inversion on
    $\oxtypjudge{\oxglobalctx}{\oxemptyctx}{\oxkontctx}{\oxvarctx}{
      \oxshift{\oxvalue}
    }{\oxsitype}{\oxvarctx^\prime}$, we know that
    $\oxtypevalidity{\oxglobalctx}{\oxemptyctx}{\oxvarctx^\prime}{\oxsitype}$.
    Since the type is valid without the frame, we know that the values
    cannot depend on that frame. Thus, we can conclude
    $\oxtypjudge{\oxglobalctx}{\oxemptyctx}{\oxkontctx}{\oxvarctx^\prime}{\oxvalue}
    {\oxsitype}{\oxvarctx^\prime}$.
  }

  %% new type is a subtype of the old type
  \oxpresline{
    \oxtunify[\oxcombine]{\oxemptyctx}{\oxkontctx}{\oxvarctx^\prime}
    {\oxsitype}{\oxsitype}{\oxvarctx^\prime}
  }{
    Immediate by \oxname{RR-Refl}.
  }

  %% new loan context is a sub-context of the old loan context
  \oxpresline{
    \exists \oxvarctx_o. \oxvarctx^\prime \oxintersect \oxvarctx_o = \oxvarctx^\prime
  }{
    $\oxvarctx_o = \oxvarctx^\prime$
  }
\end{preservation}

\oxpreservationcase{T-Framed}{\TFramed}{
  \oxframed{\oxexpr_i}
}{\EFramed}{\oxvarctx^\prime}

\begin{preservation}
  %% new store is valid
  \oxpresline{
    \oxstorevalidity{\oxglobalctx}{\oxvarctx^\prime}{\oxstore}
  }{
    Applying \Lemma{lemma:stack-envminus} to $\oxstorevalidity{\oxglobalctx}{
      \oxnewframe{\oxvarctx^\prime}{\oxframe}
    }{\oxnewframe{\oxstore}{\oxstackframe}}$ gives us
    $\oxstorevalidity{\oxglobalctx}{\oxvarctx^\prime}{\oxstore}$.
  }

  %% kontctx values are still valid
  \oxpresline{
    \oxwfkontctx{\oxglobalctx}{\oxvarctx^\prime}{\overline{\oxvalue}}{\oxkontctx}
  }{
    Inverting \oxname{WF-Temporaries} gives us $\forall i \in 1 \oxdots n. \;
    \oxtypjudge{\oxglobalctx}{\oxemptyctx}{\oxsitype_1 \oxdotsc \oxsitype_{i-1}}{\oxvarctx}{
      \oxvalue_i
    }{\oxsitype_i}{\oxvarctx}$.

    Applying \Lemma{lemma:values-dont-change} to the typing derivation (from
    \oxname{T-Framed}) for
    $\oxexpr$ (which we know is a value from \oxname{E-Framed}) gives us
    $\oxsubtypectx{\oxglobalctx}{\oxemptyctx}{\oxvarctx}{\oxvarctx^\prime}$.

    Thus, we can apply \Lemma{lemma:value-typing-related-envs} to each of the
    typing judgments to get $\forall i \in 1 \oxdots n. \;
    \oxtypjudge{\oxglobalctx}{\oxemptyctx}{\oxsitype_1 \oxdotsc \oxsitype_{i-1}}{\oxvarctx^\prime}{
      \oxvalue_i
    }{\oxsitype_i}{\oxvarctx^\prime}$.

    Finally, applying \oxname{WF-Temporaries} gives us
    $\oxwfkontctx{\oxglobalctx}{\oxvarctx^\prime}{\overline{\oxvalue}}{\oxkontctx}$.
  }

  %% new expression is well-typed
  \oxpresline{
    \oxtypjudge{\oxglobalctx}{\oxemptyctx}{\oxkontctx}{\oxvarctx^\prime}{\oxvalue}
    {\oxsitype}{\oxvarctx^\prime}
  }{
    We know from \oxname{E-Framed} that $\oxexpr$ is a value $\oxvalue$. Thus,
    we can apply \Lemma{lemma:value-typing-output} to
    $\oxtypjudge{\oxglobalctx}{\oxemptyctx}{\oxkontctx}{
      \oxnewframe{\oxvarctx}{\oxframe}
    }{\oxvalue}{\oxsitype}{
      \oxnewframe{\oxvarctx^\prime}{\oxframe^\prime}
    }$ to get $\oxtypjudge{\oxglobalctx}{\oxemptyctx}{\oxkontctx}{
      \oxnewframe{\oxvarctx^\prime}{\oxframe^\prime}
    }{\oxvalue}{\oxsitype}{
      \oxnewframe{\oxvarctx^\prime}{\oxframe^\prime}
    }$.

    We now wish to show that
    $\oxtypjudge{\oxglobalctx}{\oxemptyctx}{\oxkontctx}{\oxvarctx^\prime}{\oxvalue}
    {\oxsitype}{\oxvarctx^\prime}$. By inspecting the grammar of values and
    their typing rules, we know that the only values who depend on the
    context are pointers and closure values. But by inversion on
    $\oxtypjudge{\oxglobalctx}{\oxemptyctx}{\oxkontctx}{\oxvarctx}{
      \oxframed{\oxvalue}
    }{\oxsitype}{\oxvarctx^\prime}$, we know that
    $\oxtypevalidity{\oxglobalctx}{\oxemptyctx}{\oxvarctx^\prime}{\oxsitype}$.
    Since the type is valid without the frame, we know that the values
    cannot depend on that frame. Thus, we can conclude
    $\oxtypjudge{\oxglobalctx}{\oxemptyctx}{\oxkontctx}{\oxvarctx^\prime}{\oxvalue}
    {\oxsitype}{\oxvarctx^\prime}$.
  }

  %% new type is a subtype of the old type
  \oxpresline{
    \oxtunify[\oxcombine]{\oxemptyctx}{\oxkontctx}{\oxvarctx^\prime}
    {\oxsitype}{\oxsitype}{\oxvarctx^\prime}
  }{
    Immediate by \oxname{RR-Refl}.
  }

  %% new loan context is a sub-context of the old loan context
  \oxpresline{
    \exists \oxvarctx_o. \oxvarctx^\prime \oxintersect \oxvarctx_o = \oxvarctx^\prime
  }{
    $\oxvarctx_o = \oxvarctx^\prime$
  }
\end{preservation}

\noindent Case \oxname{T-BorrowIndex}: \\[0.5em]
\begin{tabular}{l}
  \pbox{\textwidth}{\vspace*{0.5em} From premise: \vspace*{0.5em}} \\

  \framebox[\textwidth]{
    \figuresize
    \begin{mathpar}
      \TBorrowIndex
    \end{mathpar}
  } \\

  \pbox{\textwidth}{
    \vspace*{0.5em}
    Since $\oxexpr =
    \oxref{\oxcprov}{\oxmuta}{\oxindex{\oxplaceexpr}{\oxexpr}}$, by inspection of
    the reduction rules, we know that $\oxexpr$ steps with the following rule:
    \vspace*{0.5em}
  } \\

  \framebox[\textwidth]{
    \figuresize
    \begin{mathpar}
      \inferrule[E-EvalCtx]{
        %% premises
        \oxreduce{\oxglobalctx}{\oxstore}{\oxexpr}
                 {\oxstore^\prime}{\oxexpr^\prime}
      }{
        \oxreduce{\oxglobalctx}{\oxstore}{\oxref{\oxcprov}{\oxmuta}{\oxindex{\oxplaceexpr}{\oxexpr}}}
                 {\oxstore^\prime}{\oxref{\oxcprov}{\oxmuta}{\oxindex{\oxplaceexpr}{\oxexpr^\prime}}}
      }
    \end{mathpar}
  } \\[1.5em]

\end{tabular}

We begin by applying the induction hypothesis on $\oxexpr$. We get that $\exists
\oxvarctx_i$ such that
$\oxstorevalidity{\oxglobalctx}{\oxvarctx_i}{\oxstore^\prime}$ and
$\oxwfkontctx{\oxglobalctx}{\oxvarctx_i}{\overline{\oxvalue}}{\oxkontctx}$ and
$\oxtypjudge{\oxglobalctx}{\oxemptyctx}{\oxkontctx}{\oxvarctx_i}{\oxexpr^\prime}{\oxtnum}{\oxvarctx_f^\prime}$
and
$\oxtunify[\oxcombine]{\oxemptyctx}{\oxkontctx}{\oxvarctx_f^\prime}{\oxtnum}{\oxtnum}{\oxvarctx_s}$
and there exists $\oxvarctx_o$ such that $\oxvarctx^\prime = \oxvarctx_s
\oxintersect \oxvarctx_o$. Note that since rewriting does nothing with
$\oxtnum$, $\oxvarctx_f^\prime = \oxvarctx_s$. We chose $\oxvarctx_i$ in the
lemma statement to be \framebox{$\oxvarctx_i$} and need to show:

\begin{preservation}
  %% new store is valid
  \oxpresline{
    \oxstorevalidity{\oxglobalctx}{\oxvarctx_i}{\oxstore^\prime}
  }{
    Immediate.
  }
  \oxpresline{
    \oxwfkontctx{\oxglobalctx}{\oxvarctx_i}{\overline{\oxvalue}}{\oxkontctx}
  }{
    Immediate
  }

  %% new expression is well-typed
  \oxpresline{
    \oxtypjudge{\oxglobalctx}{\oxemptyctx}{\oxkontctx}{\oxvarctx_i}{\oxref{\oxcprov}{\oxmuta}
      {\oxplaceexpr[\oxexpr^\prime]}}{\oxtref{\oxcprov}{\oxmuta}{\oxsitype}}
               {\oxvarctx_f^\prime[\oxcprov \mapsto \oxset{\oxloans^\prime}]}
  }{
    We'd like to apply \oxname{T-BorrowIndex}. In order to do so, we still need
    to show $\oxvarctx^\prime_f(\oxcprov) = \emptyset$,
    $\oxnotinclosure{\oxkontctx}{\oxvarctx^\prime_f}{\oxcprov}$,
    $\oxmusafety{\oxglobalctx}{\oxemptyctx}{\oxvarctx_f^\prime}{\oxmuta}{\oxplaceexpr}{\oxset{\oxloans^\prime}}$,
    and
    $\oxcomputetynoprov{\oxemptyctx}{\oxvarctx^\prime}{\oxmuta}{\oxplaceexpr}{\oxxitype}$.

    $\oxvarctx^\prime_f(\oxcprov) = \emptyset$ is immediate because $\oxvarctx^\prime_f(\oxcprov) \subseteq \oxvarctx^\prime(\oxcprov)$.

    Both $\oxnotinclosure{\oxkontctx}{\oxvarctx^\prime_f}{\oxcprov}$ and
    $\oxcomputetynoprov{\oxemptyctx}{\oxvarctx^\prime}{\oxmuta}{\oxplaceexpr}{\oxxitype}$
    follow from the fact that types are unchanged between $\oxvarctx^\prime$ and
    $\oxvarctx^\prime_f$.

    The last obligation is
    $\oxmusafety{\oxglobalctx}{\oxemptyctx}{\oxvarctx_f^\prime}{\oxmuta}{\oxplaceexpr}{\oxset{\oxloans^\prime}}$,
    which follows from \Lemma{lemma:ownership-safety-more-precise}
  }

  %% new type is a subtype of the old type
  \oxpresline{
    \oxtunify[\oxcombine]{\oxemptyctx}{\oxkontctx}{\oxvarctx_f^{\prime}[\oxcprov \mapsto \oxset{\oxloans^\prime}]}
    {\oxtnum}{\oxtnum}{\oxvarctx_s}
  }{
    Immediate, with $\oxvarctx_s = \oxvarctx_f^\prime[\oxcprov \mapsto \oxset{\oxloans^\prime}]$.
  }

  %% new loan context is a sub-context of the old loan context
  \oxpresline{
    \exists \oxvarctx_o. \oxvarctx_s \oxintersect \oxvarctx_o =
    \oxvarctx^\prime[\oxcprov \mapsto \oxset{\oxloans}]
  }{
    From our application of the induction hypothesis above, we have that
    $\oxvarctx^\prime = \oxvarctx_f^\prime \oxintersect \oxvarctx_o$. As we
    found in the typing judgement section above, $\oxset{\oxloans^\prime}
    \subseteq \oxset{\oxloans}$. Therefore, for this judgement, we can use
    $\oxvarctx_o[\oxcprov \mapsto \oxset{\oxloans} \setminus
      \oxset{\oxloans^\prime}]$.
  }
\end{preservation}



\noindent Case \oxname{T-Tuple}: \\[0.5em]
\begin{tabular}{l}
  \pbox{\textwidth}{\vspace*{0.5em} From premise: \vspace*{0.5em}} \\

  \framebox[\textwidth]{
    \figuresize
    \begin{mathpar}
      \TTuple
    \end{mathpar}
  } \\
  \pbox{\textwidth}{
    \vspace*{0.5em}
    Since $\oxexpr = \oxprod{\oxvalue_1 \oxdotsc \oxvalue_{i-1}, \oxexpr_i \oxdotsc \oxexpr_n}$, by inspection of the
    reduction rules, we know that $\oxexpr$ steps with the following rule:
    \vspace*{0.5em}
  } \\

  \framebox[\textwidth]{
    \figuresize
    \begin{mathpar}
      \inferrule[E-EvalCtx]{
        %% premises
        \oxreduce{\oxglobalctx}{\oxstore}{\oxexpr_i}
                 {\oxstore^\prime}{\oxexpr_i^\prime}
      }{
        \oxreduce{\oxglobalctx}{\oxstore}{\oxprod{\oxvalue_1 \oxdotsc \oxvalue_{i-1}, \oxexpr_i \oxdotsc \oxexpr_n}}
                 {\oxstore^\prime}{\oxprod{\oxvalue_1 \oxdotsc \oxvalue_{i-1}, \oxexpr_i^\prime, \oxexpr_{i+1} \oxdotsc \oxexpr_n}}
      }
    \end{mathpar}
  } \\[1.5em]
\end{tabular} \\[1em]

We want to apply our induction hypothesis to the premise of \oxname{E-EvalCtx}
with the typing judgement $\oxtypjudge{\oxglobalctx}{\oxemptyctx}{\oxkontctx,
  \oxsitype_1 \oxdotsc \oxsitype_{i-1}}{\oxvarctx_{i-1}}
{\oxexpr_i}{\oxsitype_i}{\oxvarctx_i}$. In order to do so we need to show that
$\forall \ \overline{\oxvalue}$.
$\oxwfkontctx{\oxglobalctx}{\oxvarctx_{i-1}}{\overline{\oxvalue}, \oxvalue_1
  \oxdotsc \oxvalue_{i-1}}{\oxkontctx, \oxsitype_1 \oxdotsc \oxsitype_{i-1}}$
given that
$\oxwfkontctx{\oxglobalctx}{\oxvarctx_0}{\overline{\oxvalue}}{\oxkontctx}$. So
let $\oxvalue \in \overline{\oxvalue}$. For each $\oxvalue_{j}$, where $0 < j <
i$, we can repeatedly apply \Lemma{lemma:values-dont-change} to get
$\oxsubtypectx{\oxglobalctx}{\oxemptyctx}{\oxvarctx_{j-1}}{\oxvarctx_{i-1}}$, and
then apply \Lemma{lemma:value-typing-related-envs} to get that $\forall \oxvalue_k \in \overline{\oxvalue}$. 
$\oxtypjudge{\oxglobalctx}{\oxemptyctx}{\oxkontctx, \oxsitype_1 \oxdotsc
  \oxsitype_{k-1}}{\oxvarctx_{i-1}}{\oxvalue_k}{\oxkontctx[k]}{\oxvarctx_{i-1}}$.


Now we can apply the induction hypothesis to get there is some
$\oxvarctx_{i-1}^\prime$ such that $\forall\, \overline{\oxvalue}$.
$\oxwfkontctx{\oxglobalctx}{\oxvarctx_{i-1}^\prime}{\overline{\oxvalue}}
{\oxkontctx, \oxsitype_1 \oxdotsc \oxsitype_{i-1}}$ and
$\oxtypjudge{\oxglobalctx}{\oxemptyctx}{\oxkontctx, \oxsitype_1 \oxdotsc \oxsitype_{i-1}}{\oxvarctx_{i-1}^{\prime}}
{\oxexpr_i^{\prime}}{\oxtype_i^{\textsc{si} \prime}}{\oxvarctx_i^{\prime}}$,
with $\oxstorevalidity{\oxglobalctx}{\oxvarctx_{i-1}^\prime}{\oxstore^\prime}$
and $\oxtunify[\oxcombine]{\oxemptyctx}{\oxkontctx, \oxsitype_1 \oxdotsc \oxsitype_{i-1}}{\oxvarctx_i^\prime}
{\oxtype_i^{\textsc{si}\prime}}{\oxsitype_i}{\oxvarctx_{si}}$ and there is some
$\oxvarctx_{oi}$ such that $\oxvarctx_i = \oxvarctx_{si} \oxintersect \oxvarctx_{oi}$.

We then pick $\oxvarctx_i$ in the lemma statement to be
\framebox{$\oxvarctx_{i-1}^{\prime}$} and need to show:


\begin{preservation}
  %% new store is valid
  \oxpresline{
    \oxstorevalidity{\oxglobalctx}{\oxvarctx_{i-1}^{\prime}}{\oxstore^\prime}
  }{
    Immediate.
  }
  \oxpresline{
    \oxwfkontctx{\oxglobalctx}{\oxvarctx_{i-1}^\prime}{\overline{\oxvalue}}{\oxkontctx}
  }{
    Immediate
  }

  %% new expression is well-typed
  \oxpresline{
    \begin{array}{r}
    \oxglobalctx; \, \oxemptyctx; \, \oxkontctx; \, \oxvarctx_{i-1}^{ \prime} \vdash
    \framebox{$\oxprod{\oxvalue_1 \oxdotsc \oxvalue_{i-1}, \oxexpr_i^\prime,
        \oxexpr_{i+1} \oxdotsc \oxexpr_n}$}\\
    : \oxprod{\oxsitype_1 \oxdotsc
        \oxsitype_{i-1}, \oxtype_i^{\textsc{si} \prime}, \oxsitype_{i+1}
        \oxdotsc \oxsitype_{n}} \Rightarrow \oxvarctx_n^{\prime}
    \end{array}
  }{ We want to apply \oxname{T-Tuple} in order to typecheck this tuple
    expression. We firstly want to show that we can typecheck $\oxvalue_1$
    through $\oxvalue_{i-1}$ with initial input context
    $\oxvarctx_{i-1}^{\prime}$. For each $\oxvalue_{j}$, where $0 < j < i$, we
    can repeatedly apply \Lemma{lemma:values-dont-change} to get
    $\oxsubtypectx{\oxglobalctx}{\oxemptyctx}{\oxvarctx_{j-1}}{\oxvarctx_{i-1}}$,
    and then apply \Lemma{lemma:value-typing-related-envs} to get that
    $\oxtypjudge{\oxglobalctx}{\oxemptyctx}{\oxkontctx, \oxsitype_1 \oxdotsc \oxsitype_{j-1}}{\oxvarctx_{i-1}}{\oxvalue_j}{\oxtype_j}{\oxvarctx_{i-1}}$,
    and finally applying our continuation typing hypothesis gives us
    $\oxtypjudge{\oxglobalctx}{\oxemptyctx}{\oxkontctx, \oxsitype_1 \oxdotsc \oxsitype_{j-1}}{\oxvarctx_{i-1}^\prime}{\oxvalue_j}{\oxtype_j}{\oxvarctx_{i-1}^\prime}$.

    Then we can use the result of our induction hypothesis,
    $\oxtypjudge{\oxglobalctx}{\oxemptyctx}{\oxkontctx, \oxsitype_1 \oxdotsc
      \oxsitype_{i-1}}{\oxvarctx_{i-1}^{\prime}}{\oxexpr_i^{\prime}}{\oxtype_i^{\textsc{si}
        \prime}}{\oxvarctx_i^{\prime}}$.

    And finally for $\oxexpr_{i+1} \oxdotsc \oxexpr_n$, we can repeatedly apply
    \Lemma{lemma:expressions-typing-more-precise}. }

  %% new type is a subtype of the old type
  \oxpresline{
    \oxtunify[\oxcombine]{\oxemptyctx}{\oxkontctx}{\oxvarctx_n^{\prime}}
             {\oxprod{\oxsitype_1 \oxdotsc \oxsitype_{i-1},
                 \oxtype_i^{\textsc{si} \prime}, \oxsitype_{i+1} \oxdotsc
                 \oxsitype_{n}}}{\oxprod{\oxsitype_1 \oxdotsc
                 \oxsitype_n}}{\oxvarctx_s}
  }{
    Apply \oxname{RR-Tuple} to
    \oxname{RR-Refl} for every pair of types that are identical, leaving just
    $\oxtype_i^{\textsc{si}\prime}$ and $\oxsitype_i$ for which we can use
    $\oxtunify[\oxcombine]{\oxemptyctx}
    {\oxkontctx, \oxsitype_1 \oxdotsc \oxsitype_{i-1}}{\oxvarctx_i^\prime}
    {\oxtype_i^{\textsc{si}\prime}}{\oxsitype_i}{\oxvarctx_{si}}$ from
    our induction hypothesis with
    \Lemma{lemma:rewriting-parallel-expression} and
    \Lemma{lemma:rewriting-smaller-kont} to get that
    $\oxtunify[\oxcombine]{\oxemptyctx}{\oxkontctx}{\oxvarctx_n^\prime}
    {\oxtype_i^{\textsc{si}\prime}}{\oxsitype_i}{\oxvarctx_s}$.
  }

  %% new loan context is a sub-context of the old loan context
  \oxpresline{
    \exists \oxvarctx_o. \oxvarctx_s \oxintersect \oxvarctx_o = \oxvarctx_n
  }{
    Note from the induction hypothesis above, we have $\oxvarctx_{i} =
    \oxvarctx_{si} \oxintersect \oxvarctx_{oi}$. This means that for whatever
    loan sets were unioned between $\oxvarctx_i^\prime$ and $\oxvarctx_{si}$,
    those loan sets will be unioned between $\oxvarctx^\prime_n$ and
    $\oxvarctx_s$. And since $\oxvarctx_i$ contained $\oxvarctx_{si}$, it will
    also be true that $\oxvarctx_n$ will contain $\oxvarctx_s$, which means we
    can use $\oxvarctx_o = \oxvarctx_n \setminus \oxvarctx_s$.
   }
\end{preservation}

\oxpreservationcaseheader{T-Drop}{\TDrop}

Since \oxname{T-Drop} applies to \emph{any} expression $\oxexpr$, we cannot
determine anything about what rule we stepped with. So, we will instead try to
apply our induction hypothesis to the typing derivation in the premise of
\oxname{T-Drop}
\oxtypjudge{(\oxglobalctx}{\oxemptyctx}{\oxkontctx}{
  \oxtupdate{\oxvarctx}{\oxplace}{\oxsidtype_\oxplace}
}{\oxexpr}{\oxsxtype}{\oxvarctx_f)}. To do this, we need
to establish the premises of our inductive hypothesis.

Namely, we need to show:
\begin{enumerate}
\item \oxtypjudge{\oxglobalctx}{\oxemptyctx}{\oxkontctx}{
    \oxtupdate{\oxvarctx}{\oxplace}{\oxsidtype_\oxplace}
  }{\oxexpr}{\oxsxtype}{\oxvarctx_f},
\item $\oxstorevalidity{\oxglobalctx}{
    \oxtupdate{\oxvarctx}{\oxplace}{\oxsidtype_\oxplace}
  }{\oxstore}$,
\item $\oxreduce{\oxglobalctx}{\oxstore}{\oxexpr}
  {\oxstore^\prime}{\oxexpr^\prime}$.
\end{enumerate}

\vspace{0.5em}

\begin{enumerate}
\item follows immediately from our premise.
\item follows almost directly from our premise, which tells us that
  $\oxstorevalidity{\oxglobalctx}{\oxvarctx}{\oxstore}$.
  We just need to show that the value at $\oxid$ (where $\oxid$ is the root of
  $\oxplace$) is still valid at its new type. Fortunately, it's old derivation
  works almost perfectly except for typing the part that corresponds directly to
  $\oxplace$. In this case, we can use \oxname{T-Dead} on the value to get the
  new derivation with that part of the aggregate structure at the uninitialized
  type $\oxsidtype_\oxplace$.
\item follows immediately from our premise.
\end{enumerate}

\vspace{0.5em}

This allows us to use our induction hypothesis to conclude that there exists
$\oxvarctx^\prime_i$ such that:
\begin{enumerate}
  \setcounter{enumi}{4}
\item $\oxstorevalidity{\oxglobalctx}{\oxvarctx^\prime_i}{\oxstore^\prime}$,
\item $\oxwfkontctx{\oxglobalctx}{\oxvarctx^\prime_i}{\overline{\oxvalue}}{\oxkontctx}$,
\item $\oxtypjudge{\oxglobalctx}{\oxemptyctx}{\oxkontctx}{\oxvarctx^\prime_i}
  {\oxexpr^\prime}{\oxtype^\prime}{\oxvarctx_f^\prime}$,
\item $\oxtunify[\oxcombine]{\oxemptyctx}{\oxkontctx}{\oxvarctx_f^\prime}
  {\oxtype^\prime}{\oxsxtype}{\oxvarctx_s}$, and
\item $\exists \oxvarctx_o. \; \oxvarctx_s \oxintersect \oxvarctx_o = \oxvarctx_f$.
\end{enumerate}
\hfill$\square$

%%% Local Variables:
%%% mode: latex
%%% fill-column: 80
%%% End:
